\documentclass{article}
\usepackage[margin=1in]{geometry}
\usepackage{setspace}
\usepackage{graphicx}
\usepackage{amsmath}
\usepackage{amssymb}
\usepackage{booktabs}
\usepackage{multirow}
\usepackage{tikz}
\usetikzlibrary{arrows.meta}
\usepackage[colorlinks=true,linkcolor=blue,citecolor=blue,urlcolor=blue]{hyperref}  % load last

\numberwithin{equation}{section}   % equations numbered (2.13), where "2" is the section number

\def\DDS{{\rm DDS}}
\def\DD{{\rm DD}}
\def\TDS{{\rm TDS}}
\def\LDS{{\rm LDS}}
\def\LS{{\rm LS}}
\def\LAG{{\rm LAG}}
\def\nn{\nonumber}

% These notes are split across three PDFs, compiled independently from
% notes/dedispersion.tex, notes/detrending.tex, notes/variance_map.tex.
% Everything above this comment is deliberately IDENTICAL in all three sources,
% so that a section can be moved between them without touching the preamble.
%
% A \ref cannot cross documents, so a pointer into a companion PDF names its
% section in prose instead: \ddnotes{Subband search} renders as
% [dedispersion.pdf, "Subband search"]. Each source defines the macros for the
% other two documents (either may be unused at any given moment). The argument
% must reproduce a section or subsection title of the companion verbatim.
\newcommand{\ddnotes}[1]{{\em dedispersion.pdf}, ``#1''}
\newcommand{\vmnotes}[1]{{\em variance\_map.pdf}, ``#1''}

\title{Notes on the ``pirate'' dedisperser: detrending}
\author{Kendrick Smith}

\begin{document}

\maketitle

\tableofcontents

\setstretch{1.05}   % spacing between lines
\setlength{\parskip}{2.5pt}  % spacing between paragraphs

\clearpage

\section{Time detrending algorithm 1: local polynomial subtraction}
\label{sec:time_detrending}

Upstream of the tree gridding kernel (\ddnotes{Tree gridding kernel}), each frequency channel is a 1-d intensity time
series which must be {\em detrended}: we subtract a slowly varying baseline (gain drift, sky rotation,
imperfectly calibrated bandpass) while leaving short-timescale structure -- in particular FRBs -- intact.

We have two detrenders for this problem.  The one described here is a masked, adaptively centered {\em moving
local polynomial fit}, evaluated by a van Herk block decomposition over a moment monoid.  The second, in
\S\ref{sec:kalman_detrending}, is a fixed-lag Kalman filter.  They are independent -- neither section depends
on the other -- and the intent is to run both and compare them.

\subsection{Problem statement}
\label{sec:dt_problem}

The input is a float array $d[t]$ and a boolean mask $m[t] \in \{0,1\}$, independently per (beam, frequency)
pair.  A sample is {\em valid} if $m[t]=1$ and junk (RFI-flagged, dropped packet, \ldots) if $m[t]=0$.
We model the valid samples as
\begin{equation}
d[t] = f[t] + \epsilon[t],
\end{equation}
where $f[t]$ is a smooth trend and $\epsilon[t]$ is white Gaussian noise with zero mean and $t$-independent
variance $\sigma^2$.  The detrender computes an estimate $\hat f[t]$ and replaces $d[t] \rightarrow d[t] - \hat f[t]$
at valid samples.

{\bf Outputs.}  Three per-sample arrays, of which only the first is obvious.

The detrended residual is $r[t] = d[t] - \hat f[t]$.  Alongside it the detrender emits an {\em expanded} mask
$m_{\rm out}[t]$: a sample is dropped either because its input sample was masked, or because its window turned
out to be too ill-conditioned to fit (\S\ref{sec:dt_estimator}).  Where $m_{\rm out}$ is false the residual is
meaningless and is set to zero.  So the detrender does modify the mask, and a consumer must use
$m_{\rm out}$ rather than the mask it supplied.

The third is the conditioning statistic $r_{\rm min}[t]$ of (\ref{eq:dt_rmin}), which is what the mask
expansion thresholds on.

{\bf The leverage, and why it is not one of them.}  For a {\em fixed mask} the estimator is linear in the data,
so we may write $\hat f = Hd$ for a $T \times T$ smoothing matrix $H$ (which depends on the mask but not on
$d$).  The {\em leverage} is its diagonal,
\begin{equation}
\lambda[t] \;\equiv\; H_{tt} \;=\; \frac{\partial \hat f[t]}{\partial d[t]} ,
\label{eq:dt_leverage_def}
\end{equation}
the sensitivity of the trend estimate at $t$ to the datum at $t$.  Equivalently,
\begin{equation}
{\rm Var}\big(\hat f[t]\big) = \sigma^2 \lambda[t],
\qquad
{\rm Var}\big(r[t]\big) = \sigma^2 \big(1 - \lambda[t]\big) ,
\label{eq:dt_var_r}
\end{equation}
exactly, on every surviving sample where the mask is locally constant.  So $\lambda$ is simultaneously the
variance of the trend estimate in units of $\sigma^2$, the fraction of the variance the detrender removes, and
a measure of how well determined the window is.  ${\rm Var}(r[t])$ is therefore {\em not} $t$-independent even
though ${\rm Var}(\epsilon[t])$ is: it varies with the local mask, dropping sharply near mask edges and in
sparsely populated windows.  For a fully valid window $\lambda = 9/8W = 4.4\times10^{-3}$, but under
adversarial masking $\lambda$ ranges over several decades and approaches $1$, where
${\rm Var}(r) \rightarrow 0$ and the residual carries no information at all.

None of that is computed.  It would matter if the peak-finder's normalization to sigmas
(\ddnotes{Peak-finding}) divided by $\sigma\sqrt{1-\lambda[t]}$ rather than by $\sigma$, but there is no path
by which a per-sample quantity could reach it: the peak-finding weights are built once at startup
(\vmnotes{PfAvarExact and VarianceMapExact}) from per-zone input variances supplied by the X-engine metadata.  Two things
make that acceptable rather than merely expedient.  Where the mask is locally constant $\lambda$ is the known
constant $9/8W$, i.e. a uniform $0.22\%$ rescaling of $\sigma$, degenerate with the overall normalization; and
the {\em variation} about it is confined to a small tail -- measured over $2.4\times10^5$ samples at $W=512$,
$0.07\%$ of surviving samples have $\lambda > 0.1$, which perturbs the variance of a sum over thousands of
channels by $O(10^{-4})$.  If a correction is ever wanted, the bulk of it is that single analytic constant, not
a per-sample array.

A separate question is whether to cut on $\lambda$ as well as on $r_{\rm min}$: they answer different
questions, and neither subsumes the other.  $r_{\rm min}$ asks whether the {\em fit} is well posed; $\lambda$
asks whether the {\em residual} carries information.  Measured over the same samples, every sample with
$r_{\rm min} < \epsilon$ had $\lambda \ge 0.1$, so mask expansion never discards a sample that $\lambda$ would
call healthy -- but the converse fails.  Those extra samples are well conditioned and numerically accurate, and
merely statistically degenerate: $N_v = 3$ with $n = 2$ determines a parabola exactly, so $G$ is perfectly
conditioned while $\lambda = 1$ and ${\rm Var}(r) = 0$.  That decision has been taken and is {\em not} to cut:
high leverage is not a problem for the downstream code, whereas numerical instability is, so the detrender
expands the mask on $r_{\rm min}$ only.

The detrender runs {\em incrementally}.  On the $i$-th call it is handed a buffer containing chunk $i$
(samples $iT_c \le t < (i+1)T_c$) together with $W$ samples of pre- and postpadding.  Samples $t < iT_c$ have
already been detrended and may not be modified, so each output sample is committed once and forever.
Parameters:
\begin{center}
\small
\begin{tabular}{llc}
\toprule
symbol & meaning & value \\
\midrule
$n$      & polynomial degree                       & $2$ \\
$W$      & window half-width                       & $256$ \\
$2W{+}1$ & window length (odd, symmetric)          & $513$ \\
$B = 2W$ & scan block length                       & $512$ \\
$T_c$    & time samples per chunk                  & $2048$ \\
         & buffer length $T_c + 2W$                & $2560$ \\
         & scan blocks per chunk, $T_c/B + 1$      & $5$ \\
$\epsilon$ & mask-expansion threshold on $r_{\rm min}$ & $10^{-3}$ \\
         & arithmetic                              & float32 \\
\bottomrule
\end{tabular}
\end{center}
The values are those the kernel is compiled at ({\tt include/pirate/Detrender1d.hpp}); everything below is
written for general $(n, W, T_c)$ and the numpy reference is parameterized by them.  An earlier draft of this
section used $W = 512$, and the {\em measured} numbers quoted below were taken there unless stated otherwise;
the analytic ones have been recomputed at $W = 256$.

\subsection{The estimator}
\label{sec:dt_estimator}

Fix an output time $t$ and let $S_t = \{t-W, \ldots, t+W\}$ be its window.  Define the valid count and the
{\em mask-weighted centroid}
\begin{equation}
N_v = \sum_{u \in S_t} m[u],
\qquad
c = \frac{1}{N_v}\sum_{u \in S_t} m[u]\, u,
\qquad
x_u = \frac{u - c}{W} .
\label{eq:dt_centroid}
\end{equation}
Let $\hat P$ be the degree-$n$ polynomial in $x$ minimizing $\sum_{u \in S_t} m[u]\,(d[u] - P(x_u))^2$, and set
\begin{equation}
\hat f[t] = \hat P(x_0), \qquad x_0 = -c/W,
\label{eq:dt_fhat}
\end{equation}
i.e.~the fit evaluated back at the window {\em center} $u = t$.  Both the choice of origin (the centroid) and
the choice of length unit ($W$) matter, but in opposite ways: the length unit is a diagonal rescaling of the
basis and is numerically irrelevant, whereas anchoring the expansion at the centroid rather than at the window
center is what keeps the normal equations well conditioned when the mask leaves only a narrow, off-center
sliver of valid samples.

Writing $P(x) = \sum_{j=0}^{n} a_j x^j$, the normal equations are $G a = U$ with
\begin{equation}
G_{jl} = S_{j+l},
\qquad
S_r = \sum_{u \in S_t} m[u]\, x_u^{\,r},
\qquad
U_r = \sum_{u \in S_t} (md)[u]\, x_u^{\,r},
\label{eq:dt_moments}
\end{equation}
so $G$ is the Hankel matrix of the mask moments and needs no assembly beyond indexing.
By construction $S_1 = 0$, hence $G_{01} = 0$.  For $n = 2$ this is
\begin{equation}
G = \begin{pmatrix} S_0 & 0 & S_2 \\ 0 & S_2 & S_3 \\ S_2 & S_3 & S_4 \end{pmatrix},
\qquad
U = (U_0, U_1, U_2)^T .
\label{eq:dt_G}
\end{equation}

{\bf Cholesky, and the conditioning statistic.}
$G$ is singular whenever $N_v \le n$, and ill-conditioned whenever the valid samples cannot pin down the
higher-order shape (two surviving clumps, for instance).  Rather than regularize, we detect this and drop the
sample.  Factor $G = LL^T$ in the usual way; written out for $n=2$,
\begin{align}
p_0 &= S_0,                              & L_{00} &= \sqrt{p_0}, & L_{10} &= 0, & L_{20} &= S_2/L_{00}, \nn \\
p_1 &= S_2,                              & L_{11} &= \sqrt{p_1}, & L_{21} &= S_3/L_{11}, \nn \\
p_2 &= S_4 - L_{20}^2 - L_{21}^2,        & L_{22} &= \sqrt{p_2},
\label{eq:dt_chol}
\end{align}
where in general $p_i = G_{ii} - \sum_{k<i} L_{ik}^2$.  Because $G_{01} = 0$ we get $p_0 = G_{00}$ and
$p_1 = G_{11}$ exactly, so the pivots are in natural (level, slope, curvature) order and $p_i/G_{ii} \in [0,1]$
is the fraction of order-$i$ information surviving after the lower orders are projected out.  Define
\begin{equation}
r_{\rm min}[t] \;\equiv\; \min_i \frac{p_i}{G_{ii}} , \qquad r_{\rm min} \equiv 0 \mbox{ if any } G_{ii} = 0,
\label{eq:dt_rmin}
\end{equation}
the smallest relative pivot, so that the equilibrated condition number of $G$ is roughly $1/r_{\rm min}$.  The
detrender then {\bf masks the output sample when $r_{\rm min} < \epsilon$}, with $\epsilon = 10^{-3}$.

Three remarks.  First, $r_{\rm min}$ has to look at the {\em pivots} and not at the diagonal: $N_v = 2$ with the
two samples at $\pm x$ gives $G_{ii} > 0$ for every $i$ and yet $p_2 = 0$ exactly, since $G$ has rank $2$.
Second, since a rank-$k$ $G$ has $p_i = 0$ for $i \ge k$, this masks exactly the windows that cannot determine a
degree-$n$ fit -- $N_v \le n$ always, and more besides.  It is a conditioning test and not a sample count:
$N_v = 3$ spread across the window survives, while $N_v = 300$ in two tight clumps does not.
Third, and this is the reason for choosing $r_{\rm min}$ as the threshold quantity,
\begin{equation}
r_{\rm min} \ge \epsilon
\quad\Longrightarrow\quad
{\rm cond}\big(\mbox{equilibrated } G\big) \;\le\; O(1/\epsilon) = O(10^3) ,
\label{eq:dt_cond_bound}
\end{equation}
so on every sample we keep, the relative error of the solve is bounded by
$\epsilon_{\rm mach}/\epsilon \approx 1.2\times10^{-4}$ {\em independently of the mask}.  That bound is what
makes the float32 implementation safe (\S\ref{sec:dt_numerics}).

Masking rather than regularizing costs nothing on the samples we keep, and buys exactness.  Any multiplicative
pivot floor $p_i \rightarrow \max(p_i, \epsilon G_{ii})$ would be {\em inert} wherever $r_{\rm min} \ge
\epsilon$, so the two agree on every surviving sample and differ only in the disposition of the rest.  Dropping
them instead of shrinking their fits means no surviving sample carries any shrinkage bias, so polynomial
reproduction and the identities (\ref{eq:dt_var_r}) hold exactly rather than approximately.  The price is a
mask-expansion rate measured at $5\times10^{-4}$ over the randomized-mask test suite.

One implementation detail: $p_i$ can be zero, or slightly negative through cancellation, on a sample we are
about to mask, and every sample is evaluated branch-free.  Replace $p_i$ by $\max(p_i, \mu)$ with
$\mu = 10^{-30}$ before the square root, purely to keep NaNs out of the batch.  $\mu$ is a guard, not a tuning
parameter: it is inert wherever $r_{\rm min} \ge \epsilon$.

{\bf Solve and evaluate.}  With $w = (1, x_0, \ldots, x_0^n)^T$,
\begin{equation}
\hat f[t] = w^T a, \qquad L L^T a = U ,
\end{equation}
by forward then back substitution.

{\bf Leverage.}  The leverage (\ref{eq:dt_leverage_def}) of this estimator is
\begin{equation}
\lambda[t] = \|z\|^2 = w^T G^{-1} w , \qquad Lz = w ,
\label{eq:dt_leverage}
\end{equation}
so one extra $O(n^2)$ forward substitution would give it, exactly on every surviving sample since there is no
regularizer.  It is not computed (\S\ref{sec:dt_problem}); the expression is recorded because the response
statements below are statements about it.  For a fully valid window $\lambda = 9/8W$, so the detrender removes
$0.44\%$ of the variance and ${\rm sd}(\hat f) = 1.06\,\sigma/\sqrt{W} = 0.066\,\sigma$ at $W = 256$.

{\bf Response.}  For a fully valid window the estimator is a fixed FIR smoother with
$1 - H(\omega) = O(\omega^{4})$ at low frequency and a smoothing timescale $\sim W$.  Two consequences worth
keeping in view.  A pulse of width $w \ll W$ loses a fraction $\approx 1.125\,w/W$ of its flux into the trend
($28\%$ at $w = 64$, $W = 256$), so $W$ must stay well above the widest searched boxcar $W_{\rm max}$.  And $H$ is a
truncated polynomial kernel, so $1-H$ has stopband ripple: at $n=2$ the first and largest lobe is
$H = -0.242$ at $\omega \approx 2.2\pi/W$ (exactly, $W\omega = 6.988$ in the large-$W$ limit), where the
detrender {\em amplifies} by $24\%$ rather than attenuating.  This is a property of the estimator, not of the
implementation.  Both numbers come from the closed form: the equivalent kernel is $h[s] = a + bs^2$ with
$a = 3(3P-1)/(N(4P-3))$ and $b = -15/(N(4P-3))$, where $N = 2W+1$ and $P = W(W+1)$, so
$H = a D_W + b S_W$ with $D_W$ the Dirichlet kernel $\sin(N\omega/2)/\sin(\omega/2)$ and
$S_W = -D_W''$; as $W \rightarrow \infty$ this tends to
$H(x) = \frac{15}{2}(x^{-3}\sin x - x^{-2}\cos x) - \frac{3}{2}x^{-1}\sin x$ with $x = W\omega$, which is
$1 - x^4/280 + O(x^6)$ at low frequency.  ($n=1$ is the boxcar, whose ripple is $0.217$ at
$\omega \approx 1.4\pi/W$.)

\subsection{The moment monoid}
\label{sec:dt_monoid}

Everything in (\ref{eq:dt_moments}) is computed by combining partial sums, so define for any sample set $S$
\begin{equation}
{\cal M}(S) = \big(\, N_v,\ c,\ S_2, \ldots, S_{2n},\ U_0, \ldots, U_n \,\big),
\end{equation}
each moment taken about that set's own centroid $c$, with $S_1 \equiv 0$ omitted.  This is $3n+2$ numbers
($8$ at $n=2$).  Change of origin acts by the Pascal matrix
\begin{equation}
(T_\delta v)_r = \sum_{j \le r} \binom{r}{j}\, \delta^{\,r-j}\, v_j ,
\qquad
T_\delta T_{\delta'} = T_{\delta + \delta'} ,
\label{eq:dt_shift}
\end{equation}
and for disjoint $S, S'$ the merge ${\cal M}(S) \oplus {\cal M}(S') = {\cal M}(S \cup S')$ is
\begin{equation}
N = N_v + N_v',
\quad
f = N_v'/N,
\quad
\Delta = (c' - c)/W,
\quad
c'' = c + f\Delta W,
\quad
\delta = -f\Delta,
\quad
\delta' = (1-f)\Delta ,
\label{eq:dt_merge}
\end{equation}
followed by applying $T_\delta$ and $T_{\delta'}$ to the two moment vectors and adding componentwise.
The centroid update is written in this form (rather than $(N_v c + N_v' c')/N$) to avoid forming the large
product.  Note $|\delta| + |\delta'| = |\Delta|$ exactly, so the total shift is the centroid separation and no
scheme can do better.  The $r=1$ rows need not be computed: they cancel to zero by the definition of $c''$,
which is both a flop saving and a good debug assertion.

The operation is associative and commutative (the weighted mean is associative; moments about a common origin
are set-additive), with identity $N_v = 0$, all moments zero, and $c$ arbitrary.  So $(\,\cdot\,, \oplus)$ is a
commutative monoid and supports a parallel scan.

{\bf Empty sets.}  $c$ is undefined when $N_v = 0$, and a tree scan performs many merges involving empty
aggregates (any fully masked sub-block).  Two rules make this safe and branch-free: initialize $c$ for an
empty set to a {\em finite} nominal value (its sub-block center), and guard the divide, $f := 0$ when $N = 0$.
Then merging empty with non-empty gives $f = 1$, $\delta' = 0$, $c'' = c'$, which is correct, and merging two
empties propagates zeros.  Getting this wrong injects NaNs into every window touching a masked block.

\subsection{Block decomposition and scans}
\label{sec:dt_blocks}

Partition the buffer into blocks of length $B = 2W = 512$, anchored at (chunk start $- W$); since $T_c$ is
required to be a multiple of $B$ (here $T_c = 4B$) this lattice is the same for every chunk, each chunk needs
exactly $T_c/B + 1 = 5$ blocks, and the last block of chunk $i$ is the first block of chunk $i+1$.

A window has length $B+1 = 513$, so it spans {\em exactly} two adjacent blocks for every alignment: an
interval $[q, q+B]$ meets blocks $\lfloor q/B \rfloor$ and $\lfloor q/B \rfloor + 1$ and no others.  Writing
$p = q \bmod B$,
\begin{equation}
{\cal M}(S_t) \;=\; \mathrm{Suff}_b[p] \,\oplus\, \mathrm{Pref}_{b+1}[p],
\qquad
\underbrace{(B-p)}_{\text{offsets } p \ldots B-1} + \underbrace{(p+1)}_{\text{offsets } 0 \ldots p} = B+1 ,
\label{eq:dt_vanherk}
\end{equation}
where $\mathrm{Pref}$ and $\mathrm{Suff}$ are inclusive prefix and suffix scans of the monoid within a block.
There is no special case at either end: $p=0$ takes all of block $b$ plus one sample of block $b+1$, and
$p = B-1$ the reverse.  An odd window over even blocks is what makes this work, and it also keeps the fit
exactly symmetric (so $S_{\rm odd} = 0$ exactly on a full window) and puts the evaluation point on a real
sample rather than half way between two.

{\bf The scans must be trees, and must carry the centroid.}  Two requirements, both load-bearing:
\begin{itemize}
\item Build $\mathrm{Pref}$ and $\mathrm{Suff}$ with a tree scan, not a sequential one, so that each output
depends on its leaves through $O(\log B)$ merges rather than $O(B)$.  Which tree does not matter -- the
load-bearing property is the depth, and both implementations use Hillis-Steele (the numpy reference because
work-efficiency is irrelevant there, the kernel because it is what a warp shuffle scan does) where a
work-efficient Blelloch scan would do equally well.  Measured
worst-case moment error over a full block in float32, against a float64 reference: $3.8\times10^{-5}$ for the
tree versus $4.0\times10^{-4}$ sequential, a factor of $10$.  (An earlier draft quoted $170$; that came from a
pure {\em reduction}, which returns only the root value, whereas a prefix scan must be judged by the worst of
all $B$ outputs.)  It also matches the structure of the GPU
implementation (warp-level shuffle scan, then a block-level tree), which is what makes the CPU reference
usable for validating it.
\item Every partial aggregate must carry its own centroid, per (\ref{eq:dt_merge}).  Accumulating the scans
about a fixed block center and shifting to the centroid at the end is much cheaper and completely wrong: a
prefix holding a cluster of width $w$ at distance $D$ from the block center has moments $\sim N_v (D/W)^r$,
and the final shift must cancel these down to $\sim N_v (w/W)^r$.  In the worst measured case this destroys
every digit ($100\%$ relative error in $\hat f$) where centroid-carrying scans give $5\times10^{-7}$.
\end{itemize}

{\bf Origin travel.}  Because the scans restart at every block, the accumulated origin travel along any path
is bounded by the block width, $|\Delta| \le 2$ in units of $W$, and the associated error amplification by
$(1+|\Delta|)^{2n}$.  This bound is structural -- it does not depend on tuning a re-origining interval -- and
it is the reason a sliding accumulator carried across a whole chunk is not an option (there $|\Delta|$ reaches
$T_c/W = 4$, and the amplification $1.6\times10^4$).

\subsection{Numerical requirements}
\label{sec:dt_numerics}

Everything runs in float32.  The budget follows from (\ref{eq:dt_cond_bound}): on a surviving sample the
equilibrated condition number is at most $O(1/\epsilon)$, so the relative error of the solve is at most
$\epsilon_{\rm mach}/\epsilon \approx 1.2\times10^{-4}$.  Measured worst case over the randomized-mask test
suite is $8.7\times10^{-5}\sigma$, consistent with that bound, and far below $\hat f$'s own statistical error of
$0.047\,\sigma$.  The bound holds {\em independently of the mask}, which is the property that matters: it is
what stops the accuracy from degrading as the mask distribution is made more adversarial.  Four requirements
keep us there.

\begin{itemize}
\item {\bf Subtract a constant offset first.}  Intensity data is total power, so $d[t]$ carries a large
positive mean (system temperature) while everything of interest -- the trend and the pulse alike -- lives at
the $\sigma$ level on top of it.  Every accumulator in (\ref{eq:dt_moments}) then sums terms of magnitude
$|d|$, and floating-point rounding is proportional to the accumulator's magnitude rather than to $\sigma$, so
the effective precision of $U_r$ is degraded by the factor $|d|/\sigma$.  Removing the offset costs nothing:
the polynomial basis contains the constant function, so the fit is exactly translation-equivariant, and
replacing $d$ by $d - \kappa$ shifts $\hat f$ by exactly $\kappa$ and leaves the residual $d - \hat f$
unchanged.  Nothing needs adding back, and $\kappa$ needs no accuracy at all: any value is exact, and it
affects only rounding.  What it must not be is {\em stale}.  Take $\kappa$ to be the masked mean of the buffer
in hand, not of the previous chunk: an inherited value collapses to zero whenever the previous chunk happened to
contain no valid samples, which silently disables the subtraction for the current one.  Measured, that costs
$5\times10^{-3}\sigma$ in float32 at an offset of $10^4\sigma$, against $1\times10^{-7}\sigma$ when $\kappa$
comes from the current buffer.  Using the current buffer is also cheaper (there is no state to carry) and makes
the per-chunk computation a pure function of its arguments.
The residual error scales as $3.4\,\epsilon_{\rm mach}|d-\kappa|$, so the requirement is
$|d-\kappa| \le 10^3\sigma$.  Note this bounds the {\em spread} of $d$ within one buffer rather than its
absolute level: no single additive constant helps when a large step falls inside a window, where a step of
$10^5\sigma$ still costs $2\times10^{-2}\sigma$.
Seen less numerically, the offset {\em is} the zeroth-order piece of the trend: this just splits the trend into
a crude constant handled once per chunk plus the local polynomial handled per sample.
\item {\bf Rebuild per chunk.}  Moment state is never carried across calls; each chunk's scans are built from
scratch.  (The shared block between consecutive chunks may be cached, but as a recomputation-avoidance
optimization only.)
\item {\bf Monomials in $x$, no orthogonalization.}  ${\rm cond}(G)$ looks alarming ($\sim\!10^{15}$ in raw
sample units) but is a pure diagonal-scaling artifact; Cholesky is equivariant under diagonal scaling and the
relevant quantity, the condition number of the equilibrated matrix $G_{jl}/\sqrt{G_{jj}G_{ll}}$, is $23$.
This is the standard equilibration result (van der Sluis; see Higham, {\em Accuracy and Stability of Numerical
Algorithms}): the error bound for Cholesky depends on the equilibrated condition number, not on
${\rm cond}(G)$.
Legendre or Gram polynomials would reduce that to $O(1)$, which buys nothing, and would cost the Hankel
structure of (\ref{eq:dt_moments}).
\item {\bf $\epsilon$ must clear the rounding noise in $r_{\rm min}$.}  Mask expansion is a threshold decision,
so it is only well defined if $\epsilon$ sits well above the precision to which $r_{\rm min}$ itself is
computed.  Measured worst-case disagreement between a float32 and a float64 run is $1.6\times10^{-5}$, so
$\epsilon = 10^{-3}$ has a margin of $60\times$.  Note this is the {\em only} place float32 precision feeds back
into a discrete decision; everywhere else it perturbs a continuous quantity.
\end{itemize}

Two things follow that are worth stating positively.  $r_{\rm min}$ is a ratio of a pivot to a diagonal entry,
hence covariant under diagonal rescaling of the basis, so the threshold $\epsilon$ needs no calibration against
$\sigma$, against the trend amplitude, or against the effective window width -- unlike an absolute ridge, which
has to be recalibrated as the mask narrows.  And the masking decision depends only on $G$, hence only on the
mask, so $\hat f$ remains linear in $d$ for a fixed mask; the transfer function, the identities
(\ref{eq:dt_var_r}), and Monte-Carlo noise propagation all remain valid.

\subsection{Cost}
\label{sec:dt_cost}

Per output sample, at $n=2$, counting FMAs.  A merge (\ref{eq:dt_merge}) is $\approx 60$ flops: two Pascal
shifts ($15$ for the mask part, $6$ for the data part, each), $7$ adds to combine, and $11$ for the centroid
update.  A work-efficient tree scan costs $\approx 2$ merges per element and we need both a prefix and a
suffix scan, over $2560$ input samples per $2048$ outputs:
\begin{center}
\small
\begin{tabular}{lc}
\toprule
 & flops/output \\
\midrule
prefix $+$ suffix tree scans ($4$ merges/input $\times 1.25$) & $300$ \\
window merge (\ref{eq:dt_vanherk}), one per output           & $60$ \\
Cholesky, two solves, evaluation                             & $35$ \\
\midrule
total & $\approx 400$ \\
\bottomrule
\end{tabular}
\end{center}
The tree scan roughly doubles the scan term relative to a sequential scan; we pay it for the factor of $10$
in accuracy and for GPU parity.  The cost is {\em independent of the mask}, which is the main reason for
choosing this decomposition: in a real-time system the latency budget must be met in the worst case, and RFI
is bursty, so a scheme whose cost degrades under heavy masking degrades for many channels at once.

At $n=3$ the merge grows to $\approx 115$ flops and the total to $\approx 715$, for an answer that is
{\em identical} wherever the window is fully valid (with $S_{\rm odd} = 0$ the Gram matrix is checkerboard and
$a_0$ depends only on the even block).  Degree $3$ differs from degree $2$ only where the mask is asymmetric
within the window, so we start at $n=2$ and measure whether it matters near mask edges.

One consequence of $n=3$ is worth flagging in advance, because it is structural rather than a matter of
flops: {\bf the kernel's division-free mask test does not survive it}.  At $n=2$, $G_{01} = 0$ gives
$p_0 = G_{00}$ and $p_1 = G_{11}$ exactly, so $r_0 = r_1 = 1$ and $r_{\rm min}$ collapses to the single ratio
$p_2/G_{22} = \det G/(G_{00}G_{11}G_{22})$.  That is what lets {\tt Detrender1d.cu} evaluate the fit through
the adjugate and test {\tt det >= eps*N*A*C} with no square roots.  At $n=3$ we still have $r_0 = r_1 = 1$,
but $r_2$ and $r_3$ both bind, and $\min(p_2/G_{22},\, p_3/G_{33})$ cannot be recovered from $\det G$ alone --
it needs the leading principal minors, i.e.~the factorization.  So $n=3$ means an unrolled Cholesky in the
kernel rather than an adjugate.  That costs nothing in accuracy: measured over the randomized-mask suite at
$n=2$, $W=256$, restricted to the windows the detrender actually keeps ($m[t]=1$ and $r_{\rm min} \ge
\epsilon$), the two solves agree to within a factor of two in the worst case ($2.8\times10^{-7}\sigma$ for
Cholesky against $5.1\times10^{-7}\sigma$ for the adjugate, with equal medians at $5\times10^{-9}\sigma$),
and both sit two orders of magnitude below the $8.7\times10^{-5}\sigma$ the full pipeline achieves.  The
float32 budget is spent in the moment scans, not in the solve.

Storage for the flat scheme is $3n{+}2$ floats $\times\, B \times$ two scans, i.e.~$32$ KB per block at $n=2$.
A two-level variant avoids this, and is what the kernel does: within sub-blocks of $B/32 = 16$ samples the
span is only $16/W = 0.06$ in normalized units, so no adaptive centering is needed there and the sub-block
partials are pure accumulation maintainable in registers; only the $32$ sub-block aggregates per block carry
centroids, which is $1$ KB.
Since the mask is per (beam, frequency) and there are thousands of such pairs, the natural mapping is one
thread per pair walking time sequentially, with the time axis fastest in memory.

\subsection{Validation}
\label{sec:dt_validation}

Two things make this unusually testable.  A float64, direct-summation implementation of
\S\ref{sec:dt_estimator} -- no blocks, no scans, one $O(W)$ sum per output -- is a few dozen lines and serves as
a reference for the block/scan machinery.  And the estimator has an exact fixed point: fed a degree-$\le n$
polynomial with no noise, it must return {\em identically zero} at every surviving sample.  The second is the
sharper test, because the expected answer is known analytically, with no reference implementation and no
statistical tolerance -- whatever residual appears is pure accumulated roundoff.  It is also what mask expansion
made unconditional: with a regularizer there would be a shrunk-fit carve-out to exclude.

Masks are randomized rather than enumerated.  Each row of a test mask independently draws a type (all-valid with
probability $1/2$, then Bernoulli at a random rate, long gaps, one-sided, narrow off-center clusters, two
clumps, periodic dropouts, fully masked scan blocks, sparse lattices, all-masked), randomizes that type's
parameters, and is then perturbed by a random set of fully-masked and fully-valid subintervals whose lengths
follow a power law.  Rows are independent, which is what exercises the spectator axis.

The checks worth keeping as regressions:
\begin{itemize}
\item {\bf Monoid algebra}: (\ref{eq:dt_merge}) against direct moments, associativity, commutativity, identity,
$S_1 \equiv 0$, and the empty-set rule of \S\ref{sec:dt_monoid} (the NaN trap).
\item {\bf Block decomposition}: (\ref{eq:dt_vanherk}) against direct moments for {\em every} offset $p$, which
is what pins the odd-window-over-even-block arithmetic at both endpoints.
\item {\bf Polynomial exactness}, as above, at both precisions and over several chunks.  Running it through the
chunked path is also what covers chunk stitching: a wrong block anchor or an off-by-one in the buffer slicing
shows up as a nonzero residual, and a seam shows up as structure at multiples of $T_c$.  This subsumes a
separate seam test.
\item {\bf Degenerate windows}: $N_v = 0, 1, 2$ must all be mask-expanded away, with a positive control that
$N_v = n+1$ spread across the window survives.
\item {\bf Projection identity}: build the equivalent kernel $\kappa$ by feeding unit impulses through the full
path and check $\kappa[0] = \sum_s \kappa[s]^2$.  Both sides equal $\lambda$ by (\ref{eq:dt_leverage}), but the
identity is checked from the kernel alone, so it needs nothing the solve does not already produce.  Doing this
by impulse response rather than by pushing noise through and measuring ${\rm Var}(r)$ is exact, needs no
statistical tolerance, and subsumes a Monte-Carlo variance test.
\item {\bf Kernel at zero lag}: on a full window, $\kappa[0] = (G^{-1})_{00}$ for the exact discrete Gram, which
approaches $9/8W$ as $W$ grows.  This is the sharpest check here, since the expected answer is analytic -- no
reference implementation and no statistical tolerance -- and one impulse suffices.
\item {\bf Against the float64 reference}, over the randomized masks, comparing $r$, $m_{\rm out}$ and
$r_{\rm min}$.
\item {\bf Float32 versus float64}, the whole pipeline in both.  Run the two at {\em different} $\epsilon$
($10^{-3}$ against $10^{-6}$) and compare on the intersection of their masks: with matched $\epsilon$ the two
runs make identical masking decisions, so the test never exercises a sample whose $r_{\rm min}$ roundoff was
large enough to move it across the threshold.  Also track the cumulative mask-expansion rate, which is the
diagnostic for a systematic bias in float32's $r_{\rm min}$ -- that would show up as an elevated expansion rate
rather than as a residual discrepancy.
\end{itemize}

Two properties are checked analytically rather than in code, because they are properties of the estimator and
cannot regress when the implementation changes: the $O(\omega^4)$ rolloff of $1-H$ (which is the same statement
as polynomial exactness, in the other domain) and the $1.125\,w/W$ flux loss.

\clearpage

\section{Time detrending algorithm 2: Kalman filter}
\label{sec:kalman_detrending}

This is the second detrender.  It solves the problem of \S\ref{sec:dt_problem} -- same inputs $d[t]$ and
$m[t]$, same requirement to run incrementally, same four output arrays -- and differs from
\S\ref{sec:time_detrending} in what it assumes about the trend.  There, the trend is whatever a degree-$n$
polynomial can follow across a window of $2W{+}1$ samples, and the window is what makes the estimate local.
Here the trend is a random process with a characteristic timescale, the estimate is a posterior mean, and
there is no window at all.

{\bf The structural difference in the contract.}  The estimate committed at time $t$ is
\begin{equation}
\hat f_L[t] = E\big[\, f[t] \;\big|\; d[0 .. t+L] \,\big] ,
\label{eq:kf_fixed_lag}
\end{equation}
a {\em fixed-lag} estimator: it depends on $t$ and the lag $L$ only, and not at all on where the chunk
boundaries fall.  It is therefore {\bf seam-free by construction} -- not to a small tolerance, but exactly.
(A seam is a step discontinuity in the committed $\hat f$, which is a broadband transient, i.e.~precisely
what the dedispersion search is built to find.  \S\ref{sec:time_detrending} avoids seams by making each
output depend only on a window centered on it; this section avoids them by letting the endpoint move with
$t$.)  The price is memory: the buffer handed to the $i$-th call is chunk $i$ plus $L$ samples of
postpadding and {\em no prepadding}, with $O(k^2)$ of filter state carried across the call boundary, so
chunks must be processed in order.

\subsection{The model}
\label{sec:kf_model}

Fix an integer order $k \ge 1$.  The trend is a $k$-fold integrated random walk observed in noise:
\begin{align}
\mbox{prior:} \quad & \big(\Delta^k f\big)[t] = w[t], \qquad w[t] \sim {\cal N}(0, q) \mbox{ i.i.d.}, \nn \\
\mbox{likelihood:} \quad & d[t] \sim {\cal N}\big(f[t], \sigma^2\big) \quad \mbox{at valid $t$ only},
\label{eq:kf_model}
\end{align}
where $(\Delta f)[t] = f[t+1]-f[t]$.  For $k=1$ the trend is a random walk; for $k=2$ it has a level and a
slope, and the slope wanders.

The point of (\ref{eq:kf_model}) is that it is Markov in the $k$-dimensional state
\begin{equation}
x[t] = \Big( f[t],\ (\Delta f)[t],\ \ldots,\ (\Delta^{k-1} f)[t] \Big)^T ,
\end{equation}
since $x_j[t+1] = x_j[t] + x_{j+1}[t]$ for $j < k-1$ and $x_{k-1}[t+1] = x_{k-1}[t] + w[t]$.  Writing $N$ for
the superdiagonal shift,
\begin{equation}
x[t+1] = A\,x[t] + g\,w[t], \qquad A = I + N, \qquad g = e_{k-1}, \qquad d[t] = e_0^T x[t] + \epsilon[t] .
\label{eq:kf_recursion}
\end{equation}
$A$ is the Pascal matrix, $(A^s)_{jl} = \binom{s}{l-j}$, and $A^{-1} = I - N + N^2 - \cdots$ terminates
because $N$ is nilpotent.  Both are exact integer matrices in any floating-point type.

{\bf Equivalent penalized least squares.}  With a {\em diffuse} prior on $x[0]$ -- no prior at all -- the
posterior mode of (\ref{eq:kf_model}) is the minimizer of
\begin{equation}
\chi^2(f) = \sum_t m[t]\,\big(d[t]-f[t]\big)^2 \;+\; \rho \sum_t \big(\Delta^k f\big)[t]^2 ,
\qquad \rho = \sigma^2/q ,
\label{eq:kf_chi2}
\end{equation}
whose normal equations are
\begin{equation}
{\cal N}\hat f = M d, \qquad {\cal N} \equiv M + \rho\, D_k^T D_k, \qquad M = {\rm diag}(m),
\label{eq:kf_normal}
\end{equation}
with $D_k$ the $k$-th difference matrix.  Two consequences are used throughout.

${\cal N}$ is positive definite as soon as the array holds $k$ valid samples anywhere: the null space of
$D_k$ is the polynomials of degree $<k$, and a nonzero one has at most $k-1$ roots.  So there is no
per-window rank deficiency to detect and no analogue of the ill-conditioned windows of
\S\ref{sec:dt_estimator}; the prior regularizes everything, and mask handling reduces to ``skip the
measurement update'', which is exactly correct because a masked sample carries no information.

And the estimator depends on $(k,\rho)$ only, {\em not} on $\sigma$: scaling $d \rightarrow cd$ scales
$\hat f \rightarrow c\hat f$ with $\rho$ fixed.  So $\sigma$ never has to be known or estimated, and every
formula below may be written with $\sigma^2 = 1$ and $q = 1/\rho$.  What this does assume is that $\sigma$ is
independent of $t$ within a row.

{\bf Timescales.}  Three derived quantities recur:
\begin{equation}
\tau = \rho^{1/2k}, \qquad
\ell = \frac{\tau}{\sin(\pi/2k)}, \qquad
c_k = \frac{1}{2k \sin(\pi/2k)} , \qquad \ell = 2k\,c_k\,\tau .
\label{eq:kf_scales}
\end{equation}
$\tau$ is the smoothing timescale in samples and is the parameter to think in.  $\ell$ is the correlation
length of ${\cal N}^{-1}$, whose entries decay as $\exp(-|t-t'|/\ell)$; it is what governs how much lookahead
is needed.  $c_k$ is the flux-loss constant of (\ref{eq:kf_response}).  At $k=2$, $c_2 = 0.354$ and
$\ell = \sqrt2\,\tau$.

\subsection{The estimator}
\label{sec:kf_estimator}

Because $x$ is Markov and the data split at $t$, the posterior in (\ref{eq:kf_fixed_lag}) factorizes:
\begin{equation}
p\big(x[t] \mid d[0..t+L]\big) \;\propto\;
  \underbrace{p\big(x[t] \mid d[0..t]\big)}_{\rm forward\ filter} \;\cdot\;
  \underbrace{p\big(d[t{+}1..t{+}L] \mid x[t]\big)}_{\rm backward\ likelihood} .
\label{eq:kf_split}
\end{equation}
Both factors are Gaussian in $x[t]$.  Carry each in {\em information} form -- the pair $(J,\eta)$ standing for
$\exp(-\frac12 x^TJx + \eta^Tx)$, so $J$ is a precision and $\eta = J\mu$ -- and the combination is pure
addition:
\begin{equation}
\boxed{\;J[t] = J_f[t] + J_b[t], \qquad \eta[t] = \eta_f[t] + \eta_b[t], \qquad
\hat f_L[t] = \big(J[t]^{-1}\eta[t]\big)_0 . \;}
\label{eq:kf_combine}
\end{equation}
That is the whole algorithm.  Information form rather than the more familiar covariance form for three
reasons, each of which is load-bearing: the diffuse initialization is {\em exactly} $J_f = 0$, with no
large-variance fudge, which is what makes the polynomial reproduction of \S\ref{sec:kf_exact} exact rather
than approximate; the measurement update becomes a pure addition; and the backward factor is a
{\em likelihood} rather than a distribution, so it is allowed to be rank-deficient -- which it is whenever
$[t{+}1, t{+}L]$ holds fewer than $k$ valid samples, a state the covariance form cannot represent at all.

{\bf The two recursions.}  They are mirror images: measure, rank-one downdate, then a congruence by $A^{-T}$
going forwards or $A^{T}$ going backwards.  Forwards, initialize $(J_f,\eta_f) = (0,0)$ before $t=0$ and at
each $t$
\begin{align}
\mbox{measure:}\quad & J_f \mathrel{+}= m[t]\,e_0e_0^T, \qquad \eta_f \mathrel{+}= m[t]\,d[t]\,e_0, \nn \\
\mbox{predict:}\quad & M = A^{-T}J_fA^{-1}, \quad \zeta = A^{-T}\eta_f, \quad \beta = g^TMg + 1/q, \nn \\
& J_f \leftarrow M - \frac{Mgg^TM}{\beta}, \qquad
  \eta_f \leftarrow \zeta - \frac{Mg\,(g^T\zeta)}{\beta} ,
\label{eq:kf_fwd}
\end{align}
where the post-measurement, pre-predict value is the $J_f[t]$ of (\ref{eq:kf_combine}).  Backwards, to
summarize $d[t{+}1..t{+}L]$ as a likelihood in $x[t]$, initialize $(J_b,\eta_b) = (0,0)$ at $u = t{+}L$ and
step $u$ down to $t$, absorbing the sample at $u{+}1$:
\begin{align}
& \tilde J = J_b + m[u{+}1]e_0e_0^T, \qquad \tilde\eta = \eta_b + m[u{+}1]\,d[u{+}1]\,e_0, \qquad
  \beta = g^T\tilde Jg + 1/q, \nn \\
& J_b \leftarrow A^T\Big(\tilde J - \frac{\tilde Jgg^T\tilde J}{\beta}\Big)A, \qquad
  \eta_b \leftarrow A^T\Big(\tilde\eta - \frac{\tilde Jg\,(g^T\tilde\eta)}{\beta}\Big) .
\label{eq:kf_bwd}
\end{align}
Skipping the update at $m=0$ is the entire handling of the mask.  Note the only division in either recursion
is by $\beta$, and $g^T(\cdot)g \ge 0$ by positive semidefiniteness, so $\beta \ge 1/q > 0$ always: unlike the
moment monoid of \S\ref{sec:dt_monoid} there is no guarded divide and no empty-set rule anywhere.  The
downdate is a Schur complement of a PSD matrix, so it preserves both positive semidefiniteness and rank
exactly, and ${\rm rank}\,J[t] = \min(j,k)$ where $j$ is the number of valid samples in $[0,t+L]$.

{\bf The backward window slides, and cannot be slid.}  The forward pass is a running prefix and is computed
once.  The backward factor is over $[t{+}1,t{+}L]$, which moves with $t$, and the tempting trick -- walk it by
adding the entering sample and subtracting the departing one, as \S\ref{sec:time_detrending} walks its moment
window -- does not work.  The new sample enters at the {\em far} end of the chain, which is the bottom of the
recursion; and although the single-sample step is invertible, being
$\tilde J \mapsto (\tilde J^{-1} + q\,gg^T)^{-1}$, the recursion forgets exponentially on scale $\ell$, so its
inverse amplifies exponentially.  The obstruction is exactly the rank-one downdate, i.e.~exactly the process
noise: at $q = 0$ it vanishes, the accumulation becomes additive, and $J_b$ degenerates to the masked Gram
matrix $\sum_u m[u](A^{u-t})^Te_0e_0^TA^{u-t}$ of a global polynomial fit -- which is
(\ref{eq:dt_moments}) in the binomial rather than the monomial basis, $A^s$ being the Pascal shift.  So the
moment monoid is the $q \rightarrow 0$ degeneration of this construction, and the process noise plays the
role that the finite window plays there: both exist to localize a fit that would otherwise be global.

The practical consequences are in \S\ref{sec:kf_impl}.

\subsection{Outputs}
\label{sec:kf_outputs}

The same three per-sample arrays as \S\ref{sec:dt_problem}, so the two detrenders are drop-in comparable.

{\bf Residual and expanded mask.}  Factor $J[t] = LL^T$ and, exactly as in (\ref{eq:dt_rmin}), take the
smallest relative pivot
\begin{equation}
r_{\rm min}[t] = \min_i \frac{p_i}{J_{ii}} \quad (0 \mbox{ if any } J_{ii} = 0),
\qquad
m_{\rm out}[t] = m[t] \;\wedge\; \big(r_{\rm min}[t] \ge \epsilon\big) ,
\label{eq:kf_rmin}
\end{equation}
with $\epsilon = 10^{-3}$.  This inherits the guarantee of \S\ref{sec:dt_estimator}: on a surviving sample the
equilibrated condition number is bounded, hence the relative error of the solve is
$O(\epsilon_{\rm mach}/\epsilon)$ {\em independently of the mask}.  Where $m_{\rm out}$ is false, both the
residual and $r_{\rm min}$ are set to zero, so a consumer that forgets to check the mask cannot pick up a
meaningless $r_{\rm min}$.  The expansion is a single pass: it is never fed back into $J_f$, $J_b$ or the
carried state, which is what keeps $\hat f$ linear in $d$ for a fixed input mask, and hence keeps
$m_{\rm out}$ and $r_{\rm min}$ functions of that mask alone.

Because a rank-deficient $J$ has $p_i = 0$ for $i \ge {\rm rank}$, (\ref{eq:kf_rmin}) subsumes the rank test:
``fewer than $k$ valid samples in $[0,t+L]$'' and every near-degenerate variant of it fall out of the same
comparison.  Since the fixed-lag prefix reaches back to the start of the stream, this can only fire at the start of a
stream (or after a state reset), so unlike \S\ref{sec:dt_estimator} the expansion rate is not a per-sample
property at all: it is a fixed handful of samples per stream, diluted by however long the stream runs.

{\bf Leverage.}  With $\sigma^2 = 1$,
\begin{equation}
\lambda[t] = \big(J[t]^{-1}\big)_{00} = H_{tt} = \frac{\partial \hat f_L[t]}{\partial d[t]} ,
\label{eq:kf_leverage}
\end{equation}
which the same factorization would give for free.  Like the local fit's (\ref{eq:dt_leverage}) it is not
computed (\S\ref{sec:dt_problem}); the expression is recorded because (\ref{eq:kf_response}) below is a
statement about it.  That this is the leverage and not merely the posterior variance is worth a line:
$\hat f_L[t]$ is row $t$ of the {\em full} smoother of the sub-problem on $[0,t+L]$, and for that
sub-problem the posterior precision of $f$ is ${\cal N}/\sigma^2$, so $\Sigma = \sigma^2{\cal N}^{-1}$ and
comparing with $\hat f = {\cal N}^{-1}Md$ gives $H = \Sigma M/\sigma^2$, whence $H_{tt} = \Sigma_{tt}/\sigma^2$.
The fixed-lag $H_L$ is itself {\em not} symmetric -- each row comes from a different sub-problem -- but the
diagonal identity survives because each row is a row of a symmetric matrix.  Note $0 \le \lambda \le 1$ at
every valid sample, since
$e_t^T{\cal N}^{-1}e_t = \max_v v_t^2/(v^T{\cal N}v) \le \max_v v_t^2/(v^TMv) = 1$.

{\bf Residual variance needs a second number.}  In \S\ref{sec:time_detrending} the restricted $H$ is a
projection, so ${\rm Var}(r) = \sigma^2(1-\lambda)$ exactly and one number does both jobs.  A penalized
estimator {\em shrinks} rather than projects, and instead
\begin{equation}
\frac{{\rm Var}(r[t])}{\sigma^2} = 1 - 2\lambda[t] + \nu[t] , \qquad
\nu[t] \equiv \frac{{\rm Var}(\hat f_L[t])}{\sigma^2} = \sum_u m[u]\,H_{tu}^2 \;\le\; \lambda[t] .
\label{eq:kf_var_r}
\end{equation}
Neither $\lambda$ nor $\nu$ is computed, so this matters only if the residual variance is ever wanted.  It is
recorded because it is the one place where the two detrenders genuinely differ in what a correction would
cost: $\lambda$ would be free from the factorization, whereas $\nu$ needs a Lyapunov recursion for
${\rm Cov}(\eta_f)$ and ${\rm Cov}(\eta_b)$ alongside each information recursion, roughly doubling the cost.
Using $1-\lambda$ in its place overestimates the residual variance, since $\nu \le \lambda$, and therefore
understates significance -- wrong, but in the conservative direction, and both vanish together as
$\lambda \rightarrow 1$.  How much that shortcut costs has been measured: over the randomized-mask suite
$\max(\lambda - \nu)$ ran $0.12$ to $0.19$, so on adversarial masks the discrepancy is real, while in
well-sampled stretches $\lambda \approx c_k/\tau$ and $\nu \le \lambda$ bound it below a percent.  The number
is quoted here because reproducing it means reinstating both quantities in the dense oracle of
\S\ref{sec:kf_validation}, neither of which that oracle now computes.

\subsection{Exact properties}
\label{sec:kf_exact}

Three, and they are what the implementation is tested against.

{\bf Chunk invariance, exactly.}  The forward recursion is a pure left-to-right sequence and the backward
window lies entirely inside the buffer, so the arithmetic performed for output $t$ does not depend on the
chunk decomposition at all.  Running the stream at any chunk size must agree with a single whole-stream call
{\em bit for bit}.  Moreover $J_f$ and $J_b$ are driven by the mask alone -- a measurement adds $e_0e_0^T$
regardless of $d$ -- so $m_{\rm out}$ and $r_{\rm min}$ are bit-identical unconditionally, while
the residual is bit-identical once the constant offset of \S\ref{sec:kf_impl} is held fixed.  This is the
operational meaning of ``seam-free''.

{\bf Polynomial reproduction, at two different strengths.}  Setting $f = P$ in (\ref{eq:kf_normal}) leaves the
residual $\rho\,D_k^TD_kP$, and $(D_k^TD_kP)[j] \propto (\Delta^{2k}P)[j-k]$ for $k \le j \le T'-k-1$.  Hence:
\begin{itemize}
\item {\bf Degree $\le k-1$: exact, unconditionally.}  Here $f = P$ zeroes {\em both} terms of
(\ref{eq:kf_chi2}) at once, so it is the unique minimizer for any mask, any array length and any position, and
the residual is identically zero.  Analytic, with no reference implementation and no statistical tolerance.
\item {\bf Degree $k$ to $2k-1$: exact on the infinite array.}  $\Delta^{2k}P = 0$ kills the interior rows but
not the $k$ truncated rows at each end, so reproduction fails only near an {\em array endpoint} -- and
notably {\em not} at a mask edge: inside a gap the equation reduces to $D_k^TD_kf = 0$, which $P$ satisfies,
so a degree-$(2k{-}1)$ polynomial continues straight through a gap of any length.  The fixed-lag estimator
always has an endpoint $L$ away, so in practice the residual is $O(e^{-L/\ell})$ rather than zero.  Measured
at $k=2$, $\tau=32$ on a full mask, the degree-3 residual falls
$2.5 \rightarrow 1.7\times10^{-1} \rightarrow 2.5\times10^{-2} \rightarrow 9.1\times10^{-6} \rightarrow
8.5\times10^{-13}$ as $L/\ell$ runs $2,4,8,16,32$, i.e.~as $e^{-L/\ell}$ with a slowly growing prefactor.
\item {\bf Degree $\ge 2k$: never}, by $O(1)$.
\end{itemize}

{\bf The degenerate window is well defined.}  When $[0,t+L]$ holds fewer than $k$ valid samples, $J$ is
singular; the system stays {\em consistent} nonetheless, since $\eta \in {\rm range}(J)$ by induction from
$(0,0)$, and if $m[t]=1$ then $e_0 \in {\rm range}(J)$ because the measurement at $t$ contributes $e_0e_0^T$.
So $e_0 \perp \ker J$ and $\hat f_L[t]$ is uniquely determined even though $x[t]$ is not.  With a single valid
sample $J = e_0e_0^T$ and $\eta = d[t]e_0$, giving $\hat f = d[t]$, $r = 0$, $\lambda = 1$: exactly right,
since one datum under a diffuse prior determines its own baseline.  In float32 that exact answer survives at
$k=2$, where rank deficiency forces ${\rm range}(J) = {\rm span}(e_0)$ and hence a {\em diagonal} $J$, so the
factorization is trivial and $r = 0$ comes out bitwise.  It is not relied on: (\ref{eq:kf_rmin}) cuts these
samples anyway.

\subsection{Response and numerics}
\label{sec:kf_numerics}

{\bf Transfer function.}  Where the mask is constant over a long stretch, $J_f$ reaches the fixed point of its
Riccati recursion and the estimator degenerates to a fixed IIR filter.  For the two-sided limit
$L \rightarrow \infty$, (\ref{eq:kf_normal}) diagonalizes in Fourier space and
\begin{equation}
H(\omega) = \frac{1}{1 + \rho\,\big(2\sin(\omega/2)\big)^{2k}} ,
\qquad
1 - H(\omega) = \frac{\rho\,(2\sin(\omega/2))^{2k}}{1+\rho\,(2\sin(\omega/2))^{2k}} ,
\label{eq:kf_transfer}
\end{equation}
{\em monotone} in $|\omega|$ with no stopband ripple, against the $24\%$ ripple that
\S\ref{sec:time_detrending}'s truncated polynomial kernel has near $\omega \approx 2.2\pi/W$.  Monotonicity is
immediate here, since $\sin(\omega/2)$ is monotone on $[0,\pi]$ and $z \mapsto z/(1+z)$ is increasing.  This is
the main statistical argument for the estimator.  The kernel at zero lag is
\begin{equation}
h[0] = \frac{c_k}{\tau} = \lambda_{\rm full} ,
\label{eq:kf_response}
\end{equation}
so a pulse of width $w \ll \tau$ loses a fraction $\approx c_k w/\tau$ of its flux into the trend, the
detrender removes $\approx c_k/\tau$ of the variance, and the requirement is $\tau \gg W_{\rm max}$, the
widest searched boxcar.
(\ref{eq:kf_response}) is also a closed-form check on the whole construction, and a strong one, since it
exercises the forward recursion, the backward recursion and the combine (\ref{eq:kf_combine}) together against
an analytic answer rather than against another implementation: run (\ref{eq:kf_fwd}) and (\ref{eq:kf_bwd}) to
steady state on a full mask with $L \gg \ell$, and $h[0]$ -- read off the equivalent kernel by impulses, so
that nothing beyond the shipped outputs is needed -- must approach $c_k/\tau$.  It does, from above and as
$O(\tau^{-2})$: measured relative deviations $3.0\times10^{-2}, 7.7\times10^{-3}, 2.0\times10^{-3}$ at
$\tau = 2, 4, 8$.  This is in the test suite (\S\ref{sec:kf_validation}).  Both the history behind the output
sample and the lag ahead of it must be several $\ell$, or the forward filter is still spinning up and that
dominates the deviation instead.  Finite $L$ perturbs all of this by $O(e^{-L/\ell})$; note that truncating the future
also breaks the kernel's symmetry, so $H_L$ has a small phase response where the two-sided smoother and the
local polynomial fit have none.

{\bf Conditioning.}  ${\rm cond}(J_f)$ grows as $\tau^{2(k-1)}$, because the state components $\Delta^jf$
carry scales differing by $\tau^j$ -- $1.3\times10^5$ at $k=2$, $\tau=256$.  That is a pure diagonal scaling,
under which Cholesky is equivariant (\S\ref{sec:dt_numerics}), and the quantity that matters is the
equilibrated condition number
$J_{ij}/\sqrt{J_{ii}J_{jj}}$, which is {\em flat in $\tau$}: measured $5.83$ at $k=2$ across
$\tau = 4, 16, 64, 256$, against $23$ for the local polynomial fit's Gram matrix.  So no rescaling of the
state is needed or used, and $A$, $A^{-1}$ stay exact integer matrices.  Two regimes do not spoil this.  The
matrix actually solved is better conditioned than $J_f$ alone, since combining forward and backward
information nearly symmetrizes it: at $k=2$ its equilibrated condition number tends to $1$ as $\tau$ grows
($1.016$ at $\tau=64$, $1.004$ at $\tau=256$), the steady-state pivot ratios are $(1,1)$, and
$r_{\rm min} \rightarrow 1$ -- three decades above $\epsilon$.  And iterating $J_f$
through a long masked gap with no measurements, the raw condition number runs away but the equilibrated one
{\em plateaus} by $G \approx 10\tau$ at $13.9$, with $J_{00} \sim 1/(qG^{2k-1})$ decaying only polynomially,
so it cannot underflow for any realistic gap.  (Inside such a gap (\ref{eq:kf_normal}) reduces to
$\Delta^{2k}\hat f = 0$, so the estimate is the degree-$(2k{-}1)$ discrete spline interpolating between the
surrounding data -- it degrades gracefully rather than blowing up, though its variance does grow across the
gap, which is unavoidable and physical.)

{\bf Float32.}  Everything runs in float32.  Two things are needed and both are cheap.  The constant offset of
\S\ref{sec:kf_impl}, without which the residual is a cancelling difference of two $|d|$-sized numbers; and
$\epsilon$ clearing the float32 noise in $r_{\rm min}$, since the mask expansion is the only place float32
feeds back into a discrete decision.  That noise is measured at $3\times10^{-7}$, so $\epsilon = 10^{-3}$
clears it by a factor of $3500$.  Measured worst-case float32-against-float64 residual over the
randomized-mask suite is $3\times10^{-7}\sigma$, and -- the property the carried state depends on -- it does
{\em not} grow along a stream: the forward recursion is exponentially forgetting, so roundoff injected at time
$u$ contributes $O(e^{-(t-u)/\ell})$ and the error is bounded rather than accumulating.

\subsection{Parameters}
\label{sec:kf_params}

{\bf Order.}  $k=2$.  The naive reading points the other way -- the {\em guaranteed} reproduction degrees are
$k-1 = 1$ here against $n = 2$ in \S\ref{sec:time_detrending} -- but the guaranteed degree is not what governs
the passband; the asymptotic one is, and there $k=2$ already gives $2k-1 = 3$, matching the local fit's $n+1 =
3$.  So $k=3$ would buy $O(\omega^6)$ rolloff, more than the other detrender has.  It would also cost:
the equilibrated condition number rises from $5.83$ to $35.5$, and the degenerate window of
\S\ref{sec:kf_exact} stops being benign, since for $k \ge 3$ the deficient $J$ is no longer diagonal and a
floored pivot leaks into $\hat x_0$.

{\bf Timescale and lag.}  Take $(k, \tau, L)$ as the parameters, with $\tau$ and $L$ in samples; $\rho$ and
$q$ are machine quantities with enormous dynamic range and change meaning with $k$.  To match
\S\ref{sec:time_detrending} at half-width $W$, equate the kernels at zero lag, $9/8W$ there against $c_k/\tau$
here:
\begin{equation}
\tau = \frac{8\,c_k}{9}\,W \;=\; 0.314\,W \quad (k=2) ,
\label{eq:kf_match}
\end{equation}
which also matches the flux loss and the degrees of freedom removed.  $L$ then follows from $\ell$: the
influence of data beyond $t+L$ decays as $e^{-L/\ell}$, and $L \approx 4\ell$ leaves a $\sim\!1\%$ deviation
from the ideal two-sided estimate.  At the local fit's $W = 256$ that gives $\tau = 80.5$, $\ell = 114$,
$L \approx 455$.

The lookahead is the real cost of this detrender: $L \approx 4\ell = 1.8\,W$, against exactly $W$ for
\S\ref{sec:time_detrending}, so matched on smoothing strength it roughly doubles the latency.  Against that,
it needs no prepadding at all.

\subsection{Implementation and cost}
\label{sec:kf_impl}

{\bf Chunking and state.}  On the $i$-th call the buffer is chunk $i$ plus $L$ samples of postpadding, so
$T_c + L$ samples, and the carried state is $(J_f, \eta_f)$ plus the constant offset $\kappa$ -- $O(k^2)$ per
row, five floats at $k=2$.  Outputs are committed once and never revisited.

{\bf The constant offset.}  Intensity data is total power, so $d$ carries a large positive mean while
everything of interest sits at the $\sigma$ level on top of it.  Subtract $\kappa$, the masked mean of the
buffer in hand, before filtering.  This is mathematically inert -- the basis reproduces constants, since
$\deg 0 \le k-1$, so it shifts $\hat f$ by exactly $\kappa$ and leaves the residual unchanged -- and exists
only so that the filter carries numbers of size $|d-\kappa|$.  Unlike \S\ref{sec:dt_numerics}, $\kappa$ cannot
simply be recomputed per buffer and forgotten, because the state is expressed relative to it.  Rebasing is
exact and needs no special case: shifting the data by $\Delta = \kappa'-\kappa$ shifts the level component of
the state, so $\mu \rightarrow \mu - \Delta e_0$ with $J$ untouched, i.e.
\begin{equation}
\eta_f \;\longrightarrow\; \eta_f - \Delta\,J_f e_0 .
\label{eq:kf_rebase}
\end{equation}
At stream start $J_f = \eta_f = 0$ and this is a no-op.  If a buffer has no valid samples, keep the previous
$\kappa$ rather than collapsing to zero, which would silently disable the subtraction.

{\bf Masked samples must be selected, never weighted.}  A masked sample may hold anything at all, including
NaN, and $0 \times \infty$ is NaN.  This matters more here than in \S\ref{sec:time_detrending}: a NaN reaching
the carried state destroys every subsequent output of that row, where there it would be confined to one
buffer.

{\bf Cost.}  The forward pass is $O(k^2)$ per sample and is computed once.  The backward factor is the
expensive part.  Evaluating (\ref{eq:kf_bwd}) independently for each output is $O(Lk^3)$ per sample --
$\sim\!10^4$ flops at $L = 512$, $k=2$ -- which is fine for a reference implementation but not for a kernel.
It is not the only option: the composition of these Gaussian potentials is associative, so the same van Herk
block decomposition as \S\ref{sec:dt_blocks} applies, with blocks of length $B = L$ so that
$[t{+}1,t{+}L]$ spans exactly two adjacent blocks, prefix and suffix scans within a block, and one merge per
output.  That brings the cost to $O(k^3)$ per sample, of order $600$ flops at $k=2$, comparable to the
$\approx 400$ of \S\ref{sec:dt_cost}.  The monoid element is larger than the moment monoid's $3n{+}2$: for an
interval it is the Gaussian potential relating the state at one end to the state at the other, $2k^2+3k = 14$
numbers at $k=2$.  Note also that a sequential sweep is numerically acceptable here, unlike in
\S\ref{sec:dt_blocks}: the recursion forgets exponentially, so roundoff injected at one step decays rather
than persisting, and tree depth is not load-bearing -- only parallelism is.

\subsection{Validation}
\label{sec:kf_validation}

Implemented in {\tt pirate\_frb/detrending\_1d\_kalman} and run by {\tt python -m pirate\_frb test --dt1k}.
The checks worth keeping as regressions:
\begin{itemize}
\item {\bf Against a dense oracle}: for each output, assemble (\ref{eq:kf_normal}) on $[0,t+L]$ explicitly and
solve it with dense linear algebra, sharing no code with the recursions.  Comparing the forward filter alone
and the combination (\ref{eq:kf_combine}) pins both recursions, since $J_b = J - J_f$ is then determined by
difference.  This is the test that fixes the signs and the $A$ versus $A^{-T}$ placement, and nothing else
should be trusted before it passes.
\item {\bf Polynomial reproduction}: degree $\le k-1$ identically zero under arbitrary masks, at both
precisions; and degree $2k-1$ asserted from both sides -- $O(1)$ near the stream start, suppressed as
$e^{-L/\ell}$ in the interior -- so that neither half can silently regress.
\item {\bf Chunk invariance}: several chunk sizes against a whole-stream call, requiring bit-identity for
$m_{\rm out}$ and $r_{\rm min}$ unconditionally and for the residual at fixed $\kappa$.  Writing this as a
single tolerance over all three outputs would discard the unconditional half, which is the half worth having.
\item {\bf Steady-state response}: on a full mask, the equivalent kernel at zero lag must approach
$h[0] = c_k/\tau$ as in (\ref{eq:kf_response}), from above and as $O(\tau^{-2})$.  Like the local fit's
zero-lag check this is analytic, so it needs neither the dense oracle nor a statistical tolerance, and it
exercises both recursions and the combine at once.  Note the projection identity of
\S\ref{sec:dt_validation} has no analogue here: $\sum_u m[u]H_{tu}^2 = \nu \ne \lambda$, since a penalized
estimator shrinks rather than projects.
\item {\bf Positive semidefiniteness and finiteness} over the randomized masks, including all-masked rows and
rows with a single valid sample, with $\beta \ge 1/q$ checked at every step.  The tolerance has to be
per-precision, or float32 roundoff hides a genuine float64 failure.
\item {\bf Masked data unread}: poison masked samples with NaN/Inf and require every output {\em and the
carried state} to be bit-identical.  The state half is what keeps the hazard above from being live.
\item {\bf Float32 versus float64}, which is a test of understanding rather than a debugging test: it asks
whether the error is where \S\ref{sec:kf_numerics} says it is.  It should report the fp32 noise in
$r_{\rm min}$, which is what $\epsilon$ must clear, and check that the discrepancy does not grow along a long
stream.
\end{itemize}

\clearpage

\section{2-d detrending}
\label{sec:detrending_2d}

The detrenders of \S\ref{sec:time_detrending} and \S\ref{sec:kalman_detrending} work one frequency channel at
a time, so by construction they cannot see structure {\em along} the frequency axis: bandpass ripple, coarse
channel edges, gain steps at the boundary of an RFI-flagged region.  This section describes a third detrender
which fits both axes at once -- a B-spline in frequency whose coefficients are allowed to drift as a low-order
polynomial in time -- and subtracts it.  It is implemented in {\tt pirate\_frb/detrending\_spline} and run by
{\tt python -m pirate\_frb test --dts}.

It is meant to run {\em after} one of the 1-d time detrenders and before the tree gridding kernel
(\ddnotes{Tree gridding kernel}).  The ordering is deliberate: this detrender is {\em regularized}, so unlike
\S\ref{sec:time_detrending} it carries a small shrinkage bias, and the bias is proportional to the roughness of
the baseline it is handed.  Removing the large per-channel time baselines first is therefore the cheapest way
to reduce it (\S\ref{sec:sd_omitted}).

Everything structural in this section follows from one difference: where \S\ref{sec:time_detrending}
{\em masks} the samples it cannot fit, this one {\em regularizes} them.  A masked B-spline fit has no analogue
of the moving window's ``$N_v \le n$, drop it'' rule -- a coefficient whose support happens to be entirely
flagged leaves an exactly zero row in the normal equations, and cleaning up the resulting contamination costs a
large fraction of the surviving channels.  Adding a small penalty makes the system positive definite instead,
and converts a structural failure into a bounded bias.  The rest of the design is about what that penalty must
be, what it costs, and how the residual ill-conditioning is detected.

\subsection{Problem statement}
\label{sec:sd_problem}

The input is a float array $d[f,t]$ and a boolean mask $m[f,t]$ over $0 \le f < N_{\rm freq}$, with a leading
spectator axis over beams which is not written below.  As in \S\ref{sec:dt_problem} a sample is {\em valid} if
$m=1$ and junk otherwise, and we model the valid samples as
\begin{equation}
d[f,t] = F[f,t] + \epsilon[f,t],
\end{equation}
with $\epsilon$ white and $F$ a baseline which is smooth in $f$ -- meaning representable by the spline basis
below -- and slowly varying in $t$.

{\bf The window, and why it is short.}  Fix an output time $t$ and a half-width $W$; the fit uses the $2W+1$
time samples $t+s$, $s \in \{-W,\ldots,W\}$, with the spline coefficients modelled as polynomials of degree $n$
in $s$ and evaluated back at $s = 0$.  This is the same local polynomial subtraction as
\S\ref{sec:dt_estimator}, applied to the coefficients rather than to the data.  But $W$ here is a handful of
samples, not the hundreds of \S\ref{sec:dt_problem}, because it is doing a different job.  The frequency fit is
what removes the baseline; the window exists to keep the fit from overfitting.  At $W=0$ the fit spends
$N_\phi$ degrees of freedom on $N_{\rm freq}$ channels every sample, and the leverage is $\sim N_\phi/N_{\rm
freq}$; the window divides that by $(2W+1)/(n+1)$.  Its cost is a bias if the baseline drifts faster than
degree $n$ across $2W+1$ samples, so $W$ is bounded above by the drift timescale and below by the leverage one
is willing to accept.  $(n,W) = (0,0)$ -- an independent frequency fit per time sample -- is a corner of the
same code path rather than a separate one.

{\bf Outputs.}  Three arrays, matching \S\ref{sec:dt_problem} except in shape.  The residual
$r[f,t] = d[f,t] - \hat F[f,t]$; an expanded mask $m_{\rm out}[f,t]$, zero where the residual is meaningless;
and the conditioning statistic $r_{\rm min}[t,z]$ of (\ref{eq:sd_rmin}), which is {\em per zone} rather than
per channel.  A zone (\S\ref{sec:sd_basis} below) is a contiguous block of channels whose fit is exactly
decoupled from the rest of the band, and it is the unit in which everything here is decided: the fit is
independent per zone, the conditioning statistic is per zone, and mask expansion drops a whole zone at a time.
That is a real difference from \S\ref{sec:time_detrending}, where each output sample is masked on its own
merits: here the channels of one zone at one time either all survive or all do not, because they share a
factorization.

The detrender runs incrementally, and {\bf the caller owns the padding}.  The $i$-th call is handed
$T_c + 2W$ time samples -- chunk $i$ plus $W$ of pre- and postpadding -- and emits exactly $T_c$.  Nothing is
emitted in the padding regions and no state is carried across calls, so each call is a pure function of its
arguments: there is no stream-boundary policy, no first-chunk special case, chunks may be processed in any
order, and replay is bit-reproducible.  Supplying padding that is consistent across chunks is the caller's
responsibility.

\begin{center}
\small
\begin{tabular}{lll}
\toprule
symbol & meaning & value or range \\
\midrule
$n_\phi$   & spline degree in frequency                        & $0 \ldots 3$; $2$ typical \\
$N_\phi$   & number of B-splines                               & $K + n_\phi$ if interior knots are simple \\
$K$        & number of non-empty knot intervals                & $\le 10$ tested \\
$h$        & widest knot interval, in channels                 & up to $N_{\rm freq}$ \\
$N_z$      & number of zones                                   & $\ge 1$ \\
$n$        & degree of the time polynomial                     & $0$, $1$ or $2$ \\
$W$        & window half-width, $2W{+}1 \ge n{+}1$ required    & a few samples \\
$\eta$     & regularization strength                           & $10^{-3}$ \\
$\epsilon$ & mask-expansion threshold on $r_{\rm min}$         & $3\times10^{-5}$ (fp32), $10^{-7}$ (fp64) \\
           & arithmetic                                        & float32 \\
\bottomrule
\end{tabular}
\end{center}
No production $(n_\phi, n, W, K)$ has been frozen; the code is written for general values and the test suite
sweeps $n_\phi \in \{0,1,2,3\}$, $(n,W) \in \{(0,0),(0,1),(0,3),(1,1),(1,2),(2,1),(2,2),(2,4)\}$ and
$N_{\rm freq}$ up to $3\times10^4$.  $\eta$ and $\epsilon$ are {\em not} free: they are tied to each other by a
measured scaling law (\S\ref{sec:sd_eta_eps}).

\subsection{The B-spline basis}
\label{sec:sd_basis}

A brief review, in the form the detrender uses.  Fix a degree $n_\phi \ge 0$ and a non-decreasing integer
{\em knot vector} $k_0 \le k_1 \le \cdots \le k_{N_\phi + n_\phi}$ with values in $[0, N_{\rm freq}]$.  The
$N_\phi$ basis functions are given by the Cox-de Boor recursion, starting from the indicator of a knot interval
and raising the degree one step at a time:
\begin{equation}
\phi^{(0)}_j(x) = \left\{\begin{array}{ll} 1 & k_j \le x < k_{j+1} \\ 0 & \mbox{otherwise,}\end{array}\right.
\qquad
\phi^{(r)}_j(x) = \frac{x - k_j}{k_{j+r} - k_j}\,\phi^{(r-1)}_j(x)
                + \frac{k_{j+r+1} - x}{k_{j+r+1} - k_{j+1}}\,\phi^{(r-1)}_{j+1}(x),
\label{eq:sd_coxdeboor}
\end{equation}
with the convention $0/0 \equiv 0$, and $\phi_j \equiv \phi^{(n_\phi)}_j$.  Repeated knots are fully supported
and are used heavily below, which is why the implementation evaluates (\ref{eq:sd_coxdeboor}) in its triangular
form (the NURBS book's Algorithm A2.2) restricted to the $n_\phi+1$ functions that are nonzero at $x$: in that
form no denominator can vanish at all, provided the evaluation point's own knot interval is non-empty.

Four properties are used, and only these four:
\begin{itemize}
\item {\bf Compact support.}  ${\rm supp}\,\phi_j = [k_j, k_{j+n_\phi+1})$, so at most $n_\phi+1$ basis
functions are nonzero at any $x$.  This is what makes the Gram matrix banded with half-bandwidth $n_\phi$, and
what makes the per-channel work $O(n_\phi)$ rather than $O(N_\phi)$.
\item {\bf Nonnegativity.}  $\phi_j \ge 0$.  Used in the bound (\ref{eq:sd_lammax}).
\item {\bf Smoothness by multiplicity.}  At a knot of multiplicity $\mu$ the spline is $C^{n_\phi - \mu}$.
Multiplicity is thus a design knob and not a degenerate case: it is how a known discontinuity in the bandpass
is put into the model rather than smoothed across.
\item {\bf Partition of unity.}  $\sum_j \phi_j(x) = 1$ wherever the basis is complete.  This is the one the
whole design rests on, because it says the all-ones coefficient vector {\em is} the constant function.
\end{itemize}

{\bf Conventions.}  Channel $f$ occupies $[f, f+1)$ and its datum sits at $f + \frac12$, so knot values are
channel indices, no data point ever coincides with a knot, $k_{i+1} - k_i$ is literally a channel count, and
the interval test is the ordinary half-open integer one.  The first and last knots have multiplicity exactly
$n_\phi + 1$ ({\em clamped} ends), which is what makes the basis complete -- hence a partition of unity -- on
the whole band $[0, N_{\rm freq}]$ rather than on the interior knot span only.  This is not a stylistic choice
and it does not degrade gracefully: at end multiplicity $n_\phi$ the best fit to the constant function $1$ is
already off by $0.99$ near the band edge.

{\bf Zones.}  An interior knot of multiplicity $n_\phi+1$ is a {\em zone boundary}.  No basis function
straddles one: if the boundary occupies knot indices $i \ldots i+n_\phi$ with common value $v$, then $j \le
i-1$ gives ${\rm supp}\,\phi_j \subset [\,\cdot\,, k_{i+n_\phi}) = [\,\cdot\,, v)$ and $j \ge i$ gives
${\rm supp}\,\phi_j \subset [v, \cdot\,)$.  Consequently
\begin{equation}
\mbox{$G$ and $D_1$ are {\em exactly} block diagonal across a zone boundary,}
\label{eq:sd_zone}
\end{equation}
the fits on the two sides are independent, the partition of unity holds {\em within each zone}, and the fitted
baseline is allowed to jump discontinuously at the boundary.  All three matter: the first is why flagging one
zone can never affect another, the second is the null-space statement of \S\ref{sec:sd_regulator}, and the
third is how a hard gain step is represented exactly instead of being smeared over $n_\phi+1$ knot intervals.

\subsection{The estimator}
\label{sec:sd_estimator}

Let $\psi_0, \ldots, \psi_n$ be a basis for the polynomials of degree $\le n$ on the window offsets
$s \in \{-W, \ldots, W\}$ (the specific choice is fixed in \S\ref{sec:sd_kron}, and is orthonormal).  The
model over the window is the tensor product
\begin{equation}
b(f,s) = \sum_{j=0}^{N_\phi-1} \sum_{q=0}^{n} \alpha_{jq}\, \phi_j(f)\, \psi_q(s)
       \;=\; \sum_j a_j(s)\, \phi_j(f),
\qquad a_j(s) = \sum_q \alpha_{jq} \psi_q(s),
\label{eq:sd_model}
\end{equation}
so that $b = \Phi\alpha$, where $\Phi$ is the design matrix with row index $(f,s)$, column index $(j,q)$ and
entries $\phi_j(f)\psi_q(s)$.  $\alpha$ minimizes
\begin{equation}
\chi^2(\alpha) \;=\;
\sum_{s=-W}^{W} \sum_f m[f,t+s]\,\Big(d[f,t+s] - b(f,s)\Big)^2
\;+\; \eta \sum_{s=-W}^{W} \sum_{(j,j+1)} \Big(a_{j+1}(s) - a_j(s)\Big)^2 ,
\label{eq:sd_chi2}
\end{equation}
the second sum running over pairs of adjacent coefficients {\em within a zone}.  The committed baseline is the
fit evaluated back at the window centre,
\begin{equation}
\hat F[f,t] = \sum_j a_j(0)\,\phi_j(f), \qquad a_j(0) = \sum_q \alpha_{jq}\psi_q(0) ,
\label{eq:sd_commit}
\end{equation}
which is $\alpha_{j0}$ only for a monomial basis; with an orthogonal one every even $q$ contributes.

{\bf Sufficient statistics.}  Define, at each individual time sample $u$,
\begin{equation}
G_{jl}[u] = \sum_f m[f,u]\,\phi_j(f)\phi_l(f),
\qquad
U_j[u] = \sum_f m[f,u]\,\phi_j(f)\,d[f,u] .
\label{eq:sd_GU}
\end{equation}
These are sufficient for everything the time axis does, and they are the entire interface between the two axes:
$N_\phi(n_\phi+2)$ floats per (beam, time) against $N_{\rm freq}$, a compression of about $60$ at
$N_{\rm freq} = 4096$, $K = 16$, $n_\phi = 2$.  The full-resolution arrays are therefore touched exactly twice
-- once to form $(G,U)$, once to evaluate (\ref{eq:sd_commit}) and subtract -- and everything in between runs
on the reduced arrays.

Minimizing (\ref{eq:sd_chi2}) gives the normal equations ${\cal G}\alpha = {\cal U}$ in the composite index
$I = (j,q)$:
\begin{equation}
\boxed{\;
{\cal G}_{(jq),(lr)} = \underbrace{\sum_{s=-W}^{W} \psi_q(s)\psi_r(s)\, G_{jl}[t+s]}_{\textstyle {\cal M}^{qr}_{jl}}
\;+\; \eta\,(D_1)_{jl}\,\Theta_{qr},
\qquad
{\cal U}_{(jq)} = \sum_{s=-W}^{W} \psi_q(s)\, U_j[t+s] ,
\;}
\label{eq:sd_normal}
\end{equation}
where $D_1$ is the first-difference matrix of (\ref{eq:sd_chi2}) and
\begin{equation}
\Theta_{qr} = \sum_{s=-W}^{W} \psi_q(s)\psi_r(s)
\label{eq:sd_theta}
\end{equation}
is the {\em window} Gram of the time basis.  In tensor notation the regulator is $\eta\,{\cal R}$ with
${\cal R} = D_1 \otimes \Theta$.

Four things about (\ref{eq:sd_normal}) are worth stating explicitly.

{\bf The regulator never sees the mask.}  $\Theta$ is a sum over the full window and $D_1$ depends only on the
knot vector, so ${\cal R}$ is a fixed matrix, identical at every $(f,t)$.  That is deliberate, and it is what
makes the bias predictable rather than mask-dependent; it is also what makes the exact worst-case mask
construction of \S\ref{sec:sd_validation} a linear-fractional program rather than something intractable.

{\bf Accumulate, then solve -- not solve, then smooth.}  ${\cal M}$ sums $G$ over the window {\em before}
inverting.  The distinction is invisible when nothing is masked and matters when something is, because
$\bar G^{-1}\bar U$ is not the average of the $G[t+s]^{-1}U[t+s]$, and only the former is the least-squares
estimator of (\ref{eq:sd_chi2}).  It also behaves better: a time sample whose own $G$ is rank deficient still
contributes the information it has, instead of producing a garbage coefficient that is then averaged in.

{\bf A stencil, not a van Herk decomposition.}  ${\cal M}^{qr}$ is a fixed-coefficient correlation of $G$ along
$t$, evaluated directly, at $O(W)$ per output per moment.  The monoid and block scan of
\S\ref{sec:dt_monoid}--\S\ref{sec:dt_blocks} are $O(1)$ in $W$ and are deliberately not ported: at
$W \sim 1\ldots4$ they cost the same to within a flop per full-resolution sample, and the stencil buys no block
lattice, no prefix/suffix storage, no carried partials, parallelism in $t$, and the purity of the chunk
contract.  Because the window is symmetric and the $\psi_q$ have definite parity, each stencil is either even
or odd in $s$ and the window folds in half,
\begin{equation}
\mbox{even:}\;\; \sum_{k>0} c_k\big(G[t{+}k] + G[t{-}k]\big) + c_0 G[t],
\qquad
\mbox{odd:}\;\; \sum_{k>0} c_k\big(G[t{+}k] - G[t{-}k]\big),
\label{eq:sd_parity}
\end{equation}
which halves the multiplies and, more importantly, makes the odd moments vanish {\em exactly} rather than by
cancellation of rounding errors.  The exactness is load-bearing: ``$n=1$ reduces to $n=0$ whenever the mask is
constant over the window'' depends on it.

{\bf Banded, in one index order only.}  With $I = j(n+1) + q$ -- coefficient major -- we have
$|I - I'| = |j-l|(n+1) + |q-r|$, and ${\cal M}$ vanishes for $|j-l| > n_\phi$, so ${\cal G}$ is banded with
half-bandwidth
\begin{equation}
n_b = \max(n_\phi, 1)\,(n+1) + n .
\label{eq:sd_bandwidth}
\end{equation}
The $\max$ is not cosmetic: the two contributions to (\ref{eq:sd_normal}) have {\em different} bandwidths in
$j$, $n_\phi$ for the data block and $1$ for the regulator whatever $n_\phi$ is, since $D_1$ is a difference
penalty on coefficient indices and knows nothing about the spline degree.  With the other natural order,
$I = qN_\phi + j$, the half-bandwidth is $nN_\phi + n_\phi$, which grows with $N_\phi$ and is effectively dense
(measured $2/5/8$ coefficient-major at $n = 0/1/2$ against $2/19/36$ time-major, at $n_\phi = 2$,
$N_\phi = 17$).  Since the solve is banded, the wrong order would silently turn an $O(N n_b^2)$ factorization
into an $O(N^3)$ one.  Zones stay contiguous under this ordering because $j$ is the major index, so
(\ref{eq:sd_zone}) survives it.

\subsection{The regulator}
\label{sec:sd_regulator}

The penalty in (\ref{eq:sd_chi2}) is a first difference on {\em coefficient indices} -- a P-spline penalty in
the sense of Eilers and Marx -- assembled per zone.  Three properties recommend it over the textbook
$\int (b')^2\,df$:
\begin{itemize}
\item Same null space (the constants) and, measured, the same smallest nonzero eigenvalue to $2\%$, hence the
same bias at matched regularization strength.
\item Half-bandwidth $1$ in $j$ instead of $n_\phi$.
\item It never differentiates anything, so repeated knots cannot break it.  $\int (b^{(m)})^2$ {\em diverges}
at a knot of multiplicity $\mu > n_\phi - m + 1$, and -- the real hazard -- per-interval quadrature returns a
finite matrix anyway, with a spurious extra null direction that silently destroys the positive-definiteness
guarantee below.  Since repeated interior knots are exactly how zone boundaries and reduced-continuity points
are expressed (\S\ref{sec:sd_basis}), this is not an exotic corner.
\end{itemize}
Assembly per zone is likewise not optional: a difference penalty built across a zone boundary would couple the
two sides with weight $1$, destroying both the block structure (\ref{eq:sd_zone}) and the exactness statement
below.  Since the coefficients of a zone are contiguous, ``per zone'' just means dropping the differences that
straddle a boundary.

{\bf What is reproduced exactly.}  $\Theta$ is positive definite, so $\alpha \in \ker{\cal R}$ iff
$\alpha_{jq}$ is independent of $j$ within each zone for every $q$; and by the per-zone partition of unity
(\S\ref{sec:sd_basis}) the all-ones coefficient vector of a zone is the constant function.  Hence
\begin{equation}
\ker{\cal R} \;=\; \Big\{\, b(f,s) = c_z\,P(s) \;:\; \mbox{$c_z$ constant on each zone},\; \deg P \le n \,\Big\} ,
\label{eq:sd_kernel}
\end{equation}
and any baseline of that form -- constant in frequency within each zone, times an arbitrary degree-$n$
polynomial in time -- is removed {\bf exactly, at any $\eta$, under any mask}.  This is the analogue of the
polynomial reproduction of \S\ref{sec:dt_validation}, and like it, the expected answer is analytic, which makes
it the sharpest test available.

{\bf When the system is positive definite.}  This is the statement that replaces the ``$N_v \le n$'' rank test
of \S\ref{sec:dt_estimator}, and it is a rank condition on the {\em window}, not a channel count.

\begin{quote}
{\bf Proposition.}  Zone $z$'s block of ${\cal G}$ is positive definite if and only if the zone holds at least
one valid channel at each of $\ge n+1$ {\em distinct} offsets $s$ in the window.
\end{quote}

{\em Proof.}  The data block is positive semidefinite, so a null vector of ${\cal G}$ must lie in
$\ker{\cal R}$, i.e. be of the form (\ref{eq:sd_kernel}): $b(f,s) = P(s)$ on the zone, with $\deg P \le n$.
For such a vector
\begin{equation}
\alpha^T{\cal G}\alpha = \sum_{s=-W}^{W} L_z(s)\, P(s)^2 ,
\qquad L_z(s) = \#\{f \in z : m[f,t+s] = 1\},
\end{equation}
which vanishes iff $P$ vanishes at every offset carrying data.  A nonzero polynomial of degree $\le n$ has at
most $n$ roots, so the sum is strictly positive exactly when $n+1$ offsets carry data.

So one valid channel at each of $n+1$ offsets suffices, while thousands of valid channels all at a single
offset do not.  At $(n,W) = (0,0)$ the condition reads ``the zone holds one valid channel''.  It is evaluated
structurally, by counting the offsets at which the zone's diagonal of $G$ is nonzero, rather than inferred from
the size of a pivot: the condition is exact and combinatorial, and a magnitude test would only approximate it.
A zone failing it reports $r_{\rm min} = 0$ and is masked.

{\bf The bias, which is the price.}  For noiseless data lying exactly in the span, ${\cal U}$ is the data block
applied to the true $\alpha$, so
\begin{equation}
\hat\alpha = \big(I - \eta\,{\cal G}^{-1}{\cal R}\big)\,\alpha ,
\qquad
\delta\alpha \;\equiv\; \alpha - \hat\alpha \;=\; \eta\,{\cal G}^{-1}{\cal R}\,\alpha ,
\label{eq:sd_bias}
\end{equation}
and the leftover baseline is (\ref{eq:sd_commit}) applied to $\delta\alpha$.  Two consequences.  The bias is proportional to
${\cal R}\alpha$, i.e. to the baseline's {\em roughness on the knot scale} and not to its amplitude -- which is
(\ref{eq:sd_kernel}) again, read quantitatively.  And its size is $O(\eta)$ times the amplitude: measured over
the randomized-mask suite at $\eta = 10^{-3}$, the median is $5.4\times10^{-4}$ of the baseline amplitude and
the worst seen is $6\times10^{-3}$.  $I - \eta{\cal G}^{-1}{\cal R}$ is the {\em shrinkage operator}, with
spectrum in $[0,1]$; it appears again in \S\ref{sec:sd_omitted}.

\subsection{Conditioning: equilibration, $r_{\rm min}$, and mask expansion}
\label{sec:sd_conditioning}

The solve is: equilibrate ${\cal G}$ to unit diagonal, factor by a banded Cholesky recording every pivot,
solve, unscale.  The order is the specification, not an implementation choice.

{\bf Equilibration is load-bearing twice over.}  As in \S\ref{sec:dt_numerics} it is what makes a threshold on
a pivot scale-invariant, so that $\epsilon$ needs no calibration against the units of the data.  Less obviously
-- and this is the more important half here -- it is what makes the problem well conditioned at all.  The raw
matrix ${\cal G}$ has condition number $O(h/\eta)$: a coefficient with no data under its support has diagonal
$\eta (D_1)_{jj}\Theta_{qq}$, while one sitting in the middle of a wide populated knot interval has diagonal
$O(h)$.  At $h = 3000$, $\eta = 10^{-3}$ that is already $\sim 10^5$, and it grows without bound as $\eta$
falls.  Dividing row and column by $\sqrt{{\cal G}_{II}}$ cancels {\em both} factors exactly -- the $\eta$ on
the dead coefficients and the $h$ on the live ones -- and the equilibrated condition number is $O(1)$ in both.
This is the standard van der Sluis result again, but here it is the difference between a tractable problem and
an intractable one rather than a factor of a few.

The upper end is then pinned analytically.  Writing $\hat{\cal G}$ for the equilibrated matrix,
\begin{equation}
\lambda_{\max}(\hat{\cal G}) \;\le\; \max\big(2,\; (n_\phi+1)(n+1)\big) .
\label{eq:sd_lammax}
\end{equation}
{\em Proof sketch.}  For a rank-one $xx^T$ with $\nu$ nonzero entries,
$\nu\,{\rm diag}(xx^T) - xx^T \succeq 0$.  The data block is a nonnegative combination of such terms, one per
valid $(f,s)$, each with $\nu = (n_\phi+1)(n+1)$ nonzeros by compact support.  For the regulator (with
$\Theta = I$, \S\ref{sec:sd_kron}), $2\,{\rm diag}(D_1) - D_1$ is the signless Laplacian
$\sum_{\rm edges}(e_j + e_{j+1})(e_j + e_{j+1})^T \succeq 0$.  Hence
$\max(2,(n_\phi{+}1)(n{+}1)){\rm diag}({\cal G}) - {\cal G} \succeq 0$.  Checked numerically over
random $(n_\phi, n, W, {\rm knots}, {\rm mask})$: the bound holds in every case and is attained.  So the whole
conditioning question is a {\em lower} bound on $\lambda_{\min}(\hat{\cal G})$.

{\bf The statistic.}  Let $p_I$ be the Cholesky pivots of $\hat{\cal G}$ and define, per zone,
\begin{equation}
r_{\rm min}[t,z] = \min_{I \in z}\, p_I \;\in\; [0,1],
\qquad r_{\rm min} \equiv 0 \mbox{ if the zone fails the rank test of \S\ref{sec:sd_regulator}} .
\label{eq:sd_rmin}
\end{equation}
Since $\hat{\cal G}_{II} = 1$ this is exactly the smallest {\em relative} pivot of (\ref{eq:dt_rmin}), the same
statistic the 1-d detrender thresholds on, read off a bigger matrix.
{\bf The whole zone is masked when $r_{\rm min} < \epsilon$}, and that is the only expansion rule there is.
With $\eta > 0$ there are no dead coefficients to deflate and no contamination to chase through the connected
components of $G$'s zero pattern, which an unregularized spline fit would require; and by (\ref{eq:sd_zone})
flagging one zone can never affect another.

{\bf Why $r_{\rm min}$ rather than $\lambda_{\min}$.}  They bound different errors.  Cholesky is backward
stable, so the error in the {\em coefficients} is governed by $\lambda_{\max}/\lambda_{\min}$, and by
(\ref{eq:sd_lammax}) that is $O(1/\lambda_{\min})$.  $r_{\rm min}$ is not that quantity: pivots are diagonal
entries of Schur complements, so $\lambda_{\min} \le r_{\rm min} \le 1$ always, with a gap reaching $84\times$
on real detrender matrices.  But the coefficients are not what is emitted.  The residual is evaluated at
{\em valid channels}, and for that error we have no rigorous bound in terms of either statistic -- only a
measurement, over $523$ random and adversarial configurations spanning four decades of $r_{\rm min}$, in
float32:
\begin{center}
\small
\begin{tabular}{lc}
\toprule
$r_{\rm min}$ & worst fit error \\
\midrule
$[10^{-1}, 1)$      & $1.9\times10^{-7}$ \\
$[10^{-2}, 10^{-1})$ & $1.7\times10^{-6}$ \\
$[10^{-3}, 10^{-2})$ & $1.9\times10^{-5}$ \\
below $10^{-3}$      & $1.2\times10^{-4}$ \\
\bottomrule
\end{tabular}
\end{center}
i.e.~roughly $\epsilon_{\rm mach}/(4 r_{\rm min})$.  {\em This is an empirical rule of thumb and not a bound}:
the log-log slope against $1/r_{\rm min}$ is $0.68$ rather than $1$ and the correlation is $0.62$, so ``no case
exceeded $\epsilon_{\rm mach}/r_{\rm min}$'' is an observation about the masks we generated rather than a
theorem.  A rigorous statement would need the componentwise (Skeel) bound $|{\cal G}^{-1}||{\cal G}||\alpha|$,
which is not cheap to evaluate.

With that caveat $r_{\rm min}$ is still the right choice of the two, for a structural reason rather than a
lucky one.  The small modes of $\lambda_{\min}$ are typically {\em delocalized over coefficients with no data}
-- the extremal one is a ramp across a dead run of $K$ coefficients, giving $\lambda_{\min} \sim 1/K^2$ -- and
such a mode contributes nothing at a valid channel.  $r_{\rm min}$ is a local statistic and does not see it.
Thresholding on $\lambda_{\min}$ would therefore expand the mask on zones whose fits are perfectly accurate:
measured at $N_{\rm freq} = 10^4$ with no interior knots, $n_\phi = 2$ and $\eta = 3\times10^{-3}$,
$\lambda_{\min} = 8.6\times10^{-5}$ against $r_{\rm min} = 5.1\times10^{-4}$, so the two thresholds would fire
a factor of $6$ apart on identical data.

\subsection{The time basis, and why $\epsilon$ carries over unchanged}
\label{sec:sd_kron}

$\{\psi_q\}$ is taken {\bf orthonormal} on the window -- Gram-Schmidt (in practice QR) applied to the monomials
$s^q$ -- so that $\Theta = I$ in (\ref{eq:sd_theta}).  This is the one choice in the section that moves a
threshold rather than a constant, and the reason is a small exact result.

Suppose the mask is constant across the window, $m[f,t+s]$ independent of $s$.  Then $G[t+s] = G$ and
(\ref{eq:sd_normal}) collapses to a Kronecker product,
\begin{equation}
{\cal G} = \big(G + \eta D_1\big) \otimes \Theta .
\label{eq:sd_kron}
\end{equation}
Equilibration of a Kronecker product is the Kronecker product of the equilibrations, and
$L_{A \otimes \Theta} = L_A \otimes L_\Theta$, so the relative pivots {\em multiply}:
\begin{equation}
r_{\rm min}^{\rm 2d} = r_{\rm min}^{\rm 1d} \cdot r_{\rm min}(\Theta) .
\label{eq:sd_kron_rmin}
\end{equation}
With raw monomials $r_{\rm min}(\Theta)$ is $0.18$ to $0.25$ at $n = 2$ (tending to $1 - \sqrt5/3$, worst at
small $W$), so the entire 1-d conditioning margin would be multiplied by that factor -- enough to push the
worst adversarial mask below $\epsilon$ and expand zones whose fits are perfectly accurate.  With $\Theta = I$,
$r_{\rm min}^{\rm 2d} = r_{\rm min}^{\rm 1d}$ {\em exactly}, and the threshold $\epsilon$ therefore carries
over from the pure frequency fit with no recalibration.  It costs nothing: the basis enters only through
stencil coefficients that are precomputed either way.

Two footnotes.  Orthogonalizing does forfeit the Hankel structure that monomials have -- with $\psi_q = s^q$
the block depends on $(q,r)$ only through $q+r$, so there are $2n+1$ distinct moment arrays rather than
$(n+1)(n+2)/2$ -- which is one extra stencil at $n=2$ and nothing else.  And (\ref{eq:sd_kron}) needs the mask
constant {\em across the window}, which is strictly stronger than separable: for $m[f,t] = m_f m_t$ the data
block is $G_f \otimes \Theta_{\rm valid}$ but the regulator still carries the full-window $\Theta$, so the sum
is Kronecker only when $m_t \equiv 1$.  The distinction is easy to miss and the identity fails immediately
without it.

\subsection{Choosing $\eta$ and $\epsilon$}
\label{sec:sd_eta_eps}

$\eta$ is dimensionless.  $D_1$ is defined on coefficient indices, so it is already in index coordinates and
needs no rescaling by the knot spacing; and both the shrinkage bias and the residual noise are $W$-independent,
so one value serves every window.  It wants to be {\em small}, because its cost is the undersubtraction of
(\ref{eq:sd_bias}), which \S\ref{sec:sd_omitted} argues is dangerous at any rate however rare.

What stops it is conditioning.  The extremal behaviour is governed by the {\em widest knot interval} $h$, not
by $N_{\rm freq}$ or by $K$, and follows a power law
\begin{equation}
\lambda_{\min}^{\rm worst} \;\sim\; C\,\Big(\frac{\eta}{h}\Big)^{2r/(2r+1)},
\qquad r = \min(n_\phi, 2) ,
\label{eq:sd_law}
\end{equation}
i.e.~exponent $2/3$ at $n_\phi = 1$ and $4/5$ at $n_\phi \ge 2$.  The mechanism is that the extremal
coefficient vector is one whose spline has a root of multiplicity $r$ inside the wide interval, positioned so
that the surviving data sees only the flat bottom of it; a triple root is available at $n_\phi = 3$ but does
not pay for its $D_1$ energy, which is why $r$ saturates at $2$.  Measured by exhaustive search over single
contiguous valid runs at $K = 5$, $\eta = 10^{-3}$, $h = 256 \ldots 65536$, the ratio
$\lambda_{\min}/(\eta/h)^{4/5}$ at $n_\phi = 2$ is $14.65$, $14.82$, $14.91$, $14.96$, $14.99$ -- flat to three
digits -- while the same numbers against $(\eta/h)^{2/3}$ drift by a factor of two.  The optimal run length is
$m_* \approx 0.06 \ldots 0.10\,h$.  The worst-case $r_{\rm min}$ follows the same power of $\eta$: over five
configurations from $N_{\rm freq} = 10^4$ to $3\times10^4$, cutting $\eta$ by $3$ multiplies $r_{\rm min}$ by
$0.417 \ldots 0.421$ against $3^{-4/5} = 0.415$ predicted.

So $\eta$ and $\epsilon$ move together, and are chosen jointly.  Worst-case $r_{\rm min}/\epsilon$ from an
offline exhaustive sweep:
\begin{center}
\small
\begin{tabular}{lcc}
\toprule
configuration & $(\eta,\epsilon) = (3\times10^{-3}, 10^{-4})$ & $(10^{-3}, 3\times10^{-5})$ \\
\midrule
$N_{\rm freq} = 10^4$, $K = 10$ uniform          & $23.3\times$ & $32.7\times$ \\
$N_{\rm freq} = 10^4$, no interior knots         & $5.1\times$  & $7.1\times$ \\
$N_{\rm freq} = 3\times10^4$, one $29900$-wide interval & $2.1\times$ & $3.0\times$ \\
$N_{\rm freq} = 3\times10^4$, no interior knots  & $2.1\times$  & $3.0\times$ \\
\bottomrule
\end{tabular}
\end{center}
Cutting $\eta$ from $3\times10^{-3}$ to $10^{-3}$ costs a factor $2.4$ in $r_{\rm min}$ by (\ref{eq:sd_law}),
but lets $\epsilon$ come down by $3.3$, so every margin {\em improves} by about $1.4$ while the median
shrinkage bias roughly halves.

{\bf Lowering $\epsilon$ is not free, and that is what sets the floor.}  A zone surviving at
$r_{\rm min} \sim \epsilon$ carries a float32 fit error of order $\epsilon_{\rm mach}/(4 r_{\rm min})$, and
that error is {\em coherent} -- the same dangerous category as shrinkage, not the harmless one.  Measured at
each setting's worst configuration it is $5.8\times10^{-6}$, $8.6\times10^{-6}$, $1.6\times10^{-5}$ at
$\eta = 3\times10^{-3}$, $10^{-3}$, $3\times10^{-4}$, against shrinkage biases of $1.06\times10^{-3}$,
$5.4\times10^{-4}$, $1.7\times10^{-4}$.  At $\eta = 10^{-3}$ the numerical term is still $60\times$ below the
regulator term; by $3\times10^{-4}$ the two are within a factor of $2$ of crossing.  $\eta = 3\times10^{-4}$ is
the floor -- below it, one coherent error is simply traded for another.

$\epsilon$ is a threshold on a dimensionless pivot, so its natural scale is machine epsilon rather than
anything physical; hence one value per working precision.  The float64 value is nowhere near binding
($890\times$ margin at the worst configuration above) and is left where it is.

Finally, a fact about $r_{\rm min}$ that is worth internalizing because intuition points the wrong way: {\em a
zone with a single valid channel is one of the safest masks there is}, $r_{\rm min} \approx 1.6\times10^{-2}$,
some $40\times$ better than the worst.  The null space contains the constants, so the fit simply interpolates
that one point.  The dangerous mask is a single contiguous run of a few hundred channels at a particular
fractional offset inside one wide knot interval, with everything else flagged -- a configuration nobody would
write down by hand, which has consequences for testing (\S\ref{sec:sd_validation}).

\subsection{What is deliberately not computed}
\label{sec:sd_omitted}

This subsection is the analogue of the leverage discussion in \S\ref{sec:dt_problem}, and it turns on the same
asymmetry, which is worth stating once in general form because it decides every threshold in this section.
The search sweeps $\sim$petabytes/day for ultra-rare $\sim\!10\sigma$ outliers, so two failure modes are not
comparable:
\begin{itemize}
\item {\bf Undersubtraction} -- shrinkage, or excess variance -- leaves excursions in the data and converts
noise into triggers.  With large trial factors there are many $7\sigma$ noise fluctuations, so a mechanism that
biases by $3\sigma$ with probability $10^{-5}$ can still flood the search and force the threshold up.
{\em Dangerous at any rate, however rare.}
\item {\bf Overfitting} -- leverage -- subtracts noise instead of baseline and suppresses the residual.
Converting a $20\sigma$ event to $10\sigma$ with probability $10^{-5}$ does not move the event rate.
{\em Only a significant average effect matters.}
\end{itemize}

{\bf No residual-degrees-of-freedom cut.}  For a fixed mask the estimator is linear, $\hat F = Hd$ with
$H = \Phi\,{\cal G}^{-1}\Phi^T{\rm diag}(m)$ restricted to the valid channels, and the natural summary of
overfitting is $\nu = M_z - {\rm tr}\,H$ over a zone's $M_z$ valid channels.  (Note
${\rm tr}({\cal G}^{-1}G) = {\rm tr}\,H$ exactly, by cyclicity, so the two obvious ways of writing it are the
same number.)  It would be a real cut: a zone reduced to a single valid channel is fit {\em exactly}, has
$\nu = 0$, and is nonetheless perfectly well conditioned, so $r_{\rm min}$ does not flag it.  There is
deliberately no $\nu$ cut, for three reasons.
\begin{enumerate}
\item Overfitting here has the {\em wrong sign} and cannot manufacture an excursion at all.  Restricted to the
valid samples $H$ is symmetric positive semidefinite with spectrum in $[0,1]$ -- its nonzero eigenvalues are
those of $I - \eta{\cal G}^{-1}{\cal R}$, the shrinkage operator of (\ref{eq:sd_bias}) -- so $I-H$ and hence
$(I-H)^2$ have spectrum in $[0,1]$ too, and
\begin{equation}
{\rm Var}[r] = \sigma^2\big[(I-H)^2\big]_{ff} \;\le\; \sigma^2
\end{equation}
identically.  Measured maximum $0.9968$ over $101$ masks: over-dispersion never occurs.  This is stronger than
``small on average'' -- it is ``wrong direction, always''.
\item Operationally the leverage is negligible and rigorously bounded: ${\rm tr}\,H \le N_\phi(n+1)$, so
leverage $\le N_\phi(n+1)/M_z$.  Measured at $N_{\rm freq} = 4096$: $0.0032$ fully valid, $0.0035$ at $10\%$
flagged, $0.0064$ at $50\%$ flagged, each sitting on that bound.
\item It is not cheap.  $\nu$ needs the banded entries of ${\cal G}^{-1}$ (a Takahashi recursion on the factor
already in hand), per zone per time sample, for no false-positive protection.
\end{enumerate}
One caveat is an interface concern rather than a detrender defect: at $\nu = 0$ the residual is identically
{\em zero}, not merely small, so anything downstream that estimates a variance from the data and divides by it
will meet $0/0$ on such a zone.  If that matters, the cheap fix is not $\nu$ but the count.
${\rm tr}\,H \le \min(M_z, N_\phi(n+1))$, and $\sum_{j \in z} G_{jj}$ -- already computed for the rank test of
\S\ref{sec:sd_regulator} -- lies in $[M_z/(n_\phi{+}1),\, M_z]$ by Cauchy-Schwarz and the partition of unity,
so tightening its ``$> 0$'' to ``$> c\,N_\phi(n{+}1)$'' is a complete, free and conservative surrogate.

{\bf No runtime undersubtraction statistic}, though unlike the above this is deferred rather than settled.  The
bias (\ref{eq:sd_bias}) is the failure mode that lands in the dangerous category, so it deserves more than a
scaling argument.  The strongest lever is the pipeline order already described: the 1-d time detrender has no
shrinkage, so running it first removes the large baselines before the regularized fit ever sees them, and since
the bias is proportional to roughness this should reduce the whole problem proportionally.  A runtime tripwire
does exist in principle.  From (\ref{eq:sd_normal}) there is an exact identity,
\begin{equation}
\eta\,{\cal R}\hat\alpha \;=\; \Phi^T {\rm diag}(m)\, r ,
\label{eq:sd_rho}
\end{equation}
which says the regulator term {\em is} the basis-projection of the residual -- for an unregularized fit that
projection is exactly zero by normal-equation orthogonality, so the departure from orthogonality is precisely
the undersubtracted baseline.  It costs one extra triangular solve on the factor in hand, and normalizing by
$\sum_f m_f r_f^2$ gives a dimensionless, offset-independent $\rho$.  What stops it from being a measurement is
that substituting $\hat\alpha$ for $\alpha$ multiplies it by the shrinkage operator of (\ref{eq:sd_bias}),
whose spectrum lies in $[0,1]$, so it under-reports {\em always}, and worst exactly where shrinkage is worst
(median estimate/true falls from $0.998$ to $0.346$ across four decades of true bias).  It is monotone
(log-log correlation $0.93$), hence usable as a tripwire, and the blocking question is not the statistic but the
threshold, which needs the operational distribution of $\rho$ on real data after the 1-d stage.

{\bf No offset ($\kappa$) subtraction.}  Mathematically it would be inert: a global constant is a constant on
each zone, hence in the span by the partition of unity, so the fit is translation-equivariant and the residual
is unchanged.  But it is exactly what protects float32 precision against a large DC level, for the reason
spelled out in \S\ref{sec:dt_numerics} -- intensity data is total power, and every accumulator in
(\ref{eq:sd_GU}) then sums terms of magnitude $|d|$ while the structure of interest lives at the $\sigma$
level.  Until it exists, feeding float32 data with a large offset relative to its structure loses mantissa bits
for nothing.

\subsection{Structure and cost}
\label{sec:sd_cost}

Only two passes touch the full-resolution arrays -- $(d,m) \rightarrow (G,U)$ and coefficients $\rightarrow$
residual -- and both are pure FMA against precomputed per-channel tables of $\phi_j$ and of the products
$\phi_a\phi_b$ ($9$ floats per channel at $n_\phi = 2$).  Everything else runs on the reduced arrays: the
stencil (\ref{eq:sd_parity}) at $O(W)$ per output per moment, assembly, and a banded Cholesky of size
$N = N_\phi(n+1)$ and half-bandwidth (\ref{eq:sd_bandwidth}), $O(N n_b^2)$.

Three points that are not implementation detail.

{\bf The reduction over frequency must be an explicit binary tree} over fixed-size channel blocks, not a
library sum.  Accuracy is the lesser reason (a zone can span thousands of channels).  The real one is
reproducibility: numpy's own pairwise summation blocks by {\em stride}, so its grouping over frequency changes
silently with the length of the time axis, and the same output computed in a $1$-sample chunk and a $24$-sample
chunk then differs in the last bits.  A tree over a padded channel axis depends only on the channel count, so
chunk invariance is exact.  This is a property of the algorithm, being what makes ``pure function of its
arguments'' (\S\ref{sec:sd_problem}) true bit-for-bit rather than approximately.

{\bf The cost is independent of the mask.}  Both full-resolution passes sweep every channel with no early-out
and no compaction.  Same reasoning as \S\ref{sec:dt_cost}: the real-time latency budget must be met in the
worst case, and RFI is bursty, so a scheme that degrades under heavy flagging degrades for many channels at
once.

{\bf There is no data-dependent control flow.}  Every mask-dependent decision -- the accumulation, the
equilibration guard, the pivot guards, the rank test, the residual -- is a predicated select, not a branch, so
a GPU port has no warp divergence to design around.  That is not luck: the governing rule is {\em select masked
samples away, never weight by the mask}, adopted because a flagged sample may hold anything including NaN from
a dropped packet and $0 \times {\rm NaN} = {\rm NaN}$.  The binding constraint for a kernel is instead the
working set.  Zones decouple, so the unit of work is a {\em zone} rather than the whole coefficient vector: the
banded factor, right-hand side and equilibration scales are $11 N_\phi^{\rm zone}(n{+}1)$ floats, about $800$
bytes in float32 at $n_\phi = n = 2$ with six basis functions per zone.  That is past per-thread registers --
and registers are unavailable in any case, since arbitrary runtime knots make the dimension a runtime quantity
and the loops therefore unrollable -- so one thread per (beam, zone, time sample) is workable provided the
working set lives in {\em shared} memory, which is what \texttt{Detrender2d} does.

\subsection{Validation}
\label{sec:sd_validation}

Two exact statements anchor the suite, and both are analytic, so neither needs a reference implementation or a
statistical tolerance:
\begin{itemize}
\item {\bf Exact reproduction.}  A baseline of the form (\ref{eq:sd_kernel}) -- constant in frequency within
each zone, times an arbitrary degree-$n$ polynomial in time -- must give an identically zero residual at every
surviving sample, at any $\eta$, under arbitrary masks.  Whatever appears is pure accumulated roundoff, and the
test asserts against a computed roundoff bound rather than a fixed tolerance (measured, a few percent of it).
\item {\bf The Kronecker identity} (\ref{eq:sd_kron_rmin}).  For a window-constant mask,
$r_{\rm min}^{\rm 2d} = r_{\rm min}^{\rm 1d}$ to $10^{-10}$.  This is the test that justifies carrying
$\epsilon$ over from the pure frequency fit unchanged, and it fails immediately if $\Theta \ne I$ or if the
mask family is weakened from window-constant to merely separable.
\end{itemize}

Beyond those:
\begin{itemize}
\item {\bf A dense oracle}: for each $({\rm beam}, t)$, build the design matrix $\Phi$ of (\ref{eq:sd_model})
explicitly, all $((2W{+}1)N_{\rm freq}) \times (N_\phi(n{+}1))$ of it, row-select it by the mask, and solve
with dense linear algebra.  It shares no
code with the fast path below the basis tables -- no window moments, no banded layout, no coefficient-major
ordering -- so agreement (measured $4\times10^{-15}$) is a genuine cross-check of the stencil, the assembly and
the banded factorization rather than a restatement of them.
\item {\bf The $n=1$ degeneracy, asserted from both sides.}  For a window-constant mask the $n=1$ fit must
equal the $n=0$ fit exactly (measured $2\times10^{-16}$), since the odd moments vanish by
(\ref{eq:sd_parity}); and it must {\em differ} where the mask varies within the window (measured $0.97$),
which is what keeps the first half from passing vacuously and is the only evidence that $n>0$ earns its keep.
\item {\bf The rank condition} of \S\ref{sec:sd_regulator}: a zone with data at fewer than $n+1$ distinct
offsets must be flagged, with a positive control that exactly $n+1$ offsets survives.
\item {\bf The bandwidth} (\ref{eq:sd_bandwidth}), measured from the assembled matrix.  This is a memory-safety
test rather than an accuracy one: the $\max(n_\phi,1)$ binds only at $n_\phi = 0$, where getting it wrong
writes out of bounds instead of producing a wrong answer.
\item {\bf Chunk invariance}, bit-identical across chunkings; and {\bf masked data unread}, by poisoning
flagged samples with NaN and requiring bit-identical outputs.
\item {\bf Float32 against float64}, each at its own $\epsilon$ -- deliberately differing by $10^3$, so that
the comparison is made on the intersection of two {\em different} masking decisions rather than on samples the
two runs were guaranteed to agree about.  The cumulative mask-expansion rate is tracked alongside, since a
systematic bias in float32's $r_{\rm min}$ would show up there rather than as a residual discrepancy.
\end{itemize}
Everything is parameterized over $n_\phi$ and $(n,W)$ rather than run at one configuration.  That is not
thoroughness for its own sake: $n_\phi = 0$ is where the regulator is wider than the data block, and the
bandwidth bug it exposes is invisible at every other degree.

{\bf Generating masks, which is most of the work.}  The central fact is that {\em random masks do not find the
hard cases}.  Measured at $n_\phi = 2$, $K = 5$, $h = 3000$, $\eta = 10^{-3}$: five thousand Bernoulli masks at
log-uniform density reach a worst $r_{\rm min}$ of $9.6\times10^{-3}$, while a constructed adversarial mask
reaches $4.1\times10^{-4}$ -- a factor of $23$.  A suite built on Bernoulli masks would report an enormous
margin over $\epsilon$ no matter how wrong the conditioning handling was.

The worst case can be constructed exactly, because of a structural accident worth recording.  Fix a coefficient
vector $\alpha$ and let $w_{fs} \in \{0,1\}$ be the mask.  The equilibrated Rayleigh quotient is
{\em linear-fractional} in $w$:
\begin{equation}
R(w) = \frac{\sum_{f,s} w_{fs}\,a_{fs} + A_0}{\sum_{f,s} w_{fs}\,c_{fs} + B_0},
\qquad
a_{fs} = b(f,s)^2,
\qquad
c_{fs} = \sum_{jq} \alpha_{jq}^2\,\phi_j(f)^2\psi_q(s)^2 ,
\label{eq:sd_dinkelbach}
\end{equation}
with $A_0 = \eta\,\alpha^T{\cal R}\alpha$ and $B_0 = \eta\sum_{jq}\alpha_{jq}^2 (D_1)_{jj}\Theta_{qq}$
{\em mask-independent} -- precisely because the regulator never sees the mask (\S\ref{sec:sd_estimator}).  So
at level $\tau$ the minimizing mask is exactly $\{(f,s) : a_{fs} < \tau c_{fs}\}$, and Dinkelbach's algorithm
converges to the exact minimizer of $R$ over all $2^{N_{\rm freq}(2W+1)}$ masks at fixed $\alpha$; alternating
that with $\alpha \leftarrow$ (bottom eigenvector) gives the adversarial mask.  Verified two ways: the
decomposition reproduces the true Rayleigh quotient to $6\times10^{-16}$, and no single $(f,s)$ flip improves
the converged mask in $40$ of $40$ configurations.

The generator is then a mixture, drawn independently per zone: the adversarial construction; structured
families taken from the measured extremal geometry and from the code's corner cases (dead zones, leading versus
trailing dead runs, $M_z = 0, 1, 2$, combs, whole knot blocks); about $10\%$ separable products of two 1-d
masks; and plain Bernoulli for breadth.  The result is then perturbed by a random set of fully-flagged and
fully-valid rectangles whose extents follow a power law, which is what finds problems nobody predicted.  The
knot vector is randomized too and targets $h$ {\em directly} rather than hoping a uniform draw produces a wide
interval -- including the $K=1$ case where the whole band is a single knot interval, which by
(\ref{eq:sd_law}) is the extreme.

\clearpage

\end{document}
